% Upstream source: ~/mm/papers/speconsense/speconsense-0.8-intro.tex
% This copy is distributed alongside speconsense for convenient access; edits
% should be made upstream and synced here, not the other way around.
\documentclass[11pt]{article}

\usepackage[margin=1in]{geometry}
\usepackage{amsmath,amssymb}
\usepackage{graphicx}
\usepackage{booktabs}
\usepackage{hyperref}
\usepackage{listings}
\usepackage{xcolor}
\usepackage[numbers]{natbib}
\usepackage{microtype}
\usepackage{lineno}

\lstset{
  basicstyle=\ttfamily\small,
  breaklines=true,
  frame=single,
  framesep=4pt,
  backgroundcolor=\color{gray!5},
  showstringspaces=false,
}

\title{Speconsense: a brief history and a guided tour of the 0.8.0 pipeline}
\author{Josh Walker}
\date{Internal review draft, May 2026 (commit 896ed16)}

\begin{document}
\sloppy
\maketitle
\linenumbers

\begin{abstract}
Speconsense is a per-specimen consensus tool for Oxford Nanopore amplicon
reads, developed for the fungal DNA barcoding workflow that previously used
NGSpeciesID~\cite{ngspeciesid}, with automatic variant phasing as a
first-class feature rather than a post-hoc step. The tool has been shaped over roughly seven months
of production use at Mycota Lab (since November 2025), and the 0.8.0 release
reorganizes the read-to-consensus path into thirteen sequential phases (six
new in 0.8.0) that address three persistent concerns from the validator
community: aggressive
homopolymer normalization that hides real variation at homopolymer interfaces,
ambiguity-code inflation in low-quality clusters, and the lack of a defensible
test for low-support variants. Against the validator-facing pass track on a
955-specimen production run, 0.8.0 produces 27\% more distinct candidate
biological variants with 44\% fewer ambiguity codes and a 5\% yield gain
over 0.7.
Pass-track consensuses are also cleaner at the homopolymer interfaces
where 0.7 most often hid real variation: the share of pass-track
clusters with unphased homopolymer-interface heterogeneity at
length-2-to-5 HP boundaries drops 8-fold (27\% $\to$ 3\%). Variants the
new significance test flags as statistically indistinguishable from
basecaller noise are routed to a separate inspection track rather than
emitted to the main output. The routing thresholds are deliberate first
cuts; an 0.8.1 fast-follow tightens the MAD outlier-removal default
(\texttt{--mad-z-threshold} 3.0~$\to$~1.5; §\ref{sec:looking-ahead}),
and a 0.8.2 revision will revisit the size-based thresholds based on
partner review and user-acceptance testing. This
paper is a guided tour of the pipeline as it stands, told through the
development decisions that produced it, with empirical results from a per-phase
additive ladder and cross-tool comparisons. The companion technical papers --- variant
significance~\cite{cer_paper}, homopolymer error rate~\cite{hp_paper}, and
CER in practice~\cite{cer_in_practice} --- provide the mathematical foundations.
\end{abstract}

\section{Origins}

Speconsense exists because a typical fungal ITS specimen is more interesting
than a single sequence. A real specimen often contains genuine intragenomic
variants, occasional contaminants, and the usual nanopore noise floor --- and
a barcoding pipeline that emits one consensus per specimen will silently
average across all of them. The Russell barcoding
protocol~\cite{russell2023} used NGSpeciesID~\cite{ngspeciesid}, which clusters
reads but stops short of phasing variants and can fold biologically distinct
sequences into a single output when their cluster shapes happen to overlap.

Initial development began in February 2025 with two concrete goals: (1) cluster
reads in a way that could discover variants without being told how many to
expect, and (2) generate a separate consensus for each variant, with all the
bookkeeping needed to trace each output back to its supporting reads. After a
small handful of commits in February and March, the project went on hiatus for
roughly five months. It resumed in late August 2025 when Harte Singer at the
Fungal Diversity Survey (FUNDIS) raised the need for SNP calling at scale
across a large fungal-barcoding dataset; the author proposed speconsense as a
substitute for the NGSpeciesID-based step in the existing workflow, and serious
development began that week. A guiding principle, articulated explicitly later,
was \emph{split rather than conflate} --- when in doubt, emit two clusters and
let the post-processing tool merge them, rather than collapse them in core and
lose the distinction forever.

\section{Choosing MCL}

The dominant question for the clustering step was: how do we cluster reads
when we don't know in advance how many clusters there should be? K-means and
similar partition-based methods require the cluster count up front.
Hierarchical clustering produces a dendrogram but pushes the cut-height
decision onto the user. Greedy clustering as a class spans many
implementations --- some, like \texttt{vsearch --cluster\_fast}, use a
k-mer index and scale gracefully. The straightforward greedy variant
that speconsense initially considered, however, tends to produce a
single dominant cluster plus a long tail of fragments, and (as we will
see in Section~\ref{sec:graph}) requires a dense pairwise comparison
matrix that becomes prohibitive at scale.

Markov Clustering (MCL)~\cite{vandongen2000} was chosen as the default because
it sits in the right part of the design space: it operates on a similarity
graph, it discovers the cluster count from the graph's structure rather than as
a parameter, and its inflation parameter gives a single intuitive knob for
cluster granularity. Importantly, MCL operates on a \emph{sparse} graph ---
typically only the strongest edges from each read are retained --- which
scales gracefully to tens of thousands of reads. A greedy algorithm
(speconsense's own implementation, working from the same dense pairwise
matrix MCL builds) is retained as a fallback when MCL is not installed,
and as an option for users whose primary use case is contaminant detection
rather than variant resolution.

\subsection{Building the similarity graph}
\label{sec:graph}

MCL needs a graph as input, and the graph is what determines the clustering.
Speconsense builds it by comparing every read against every other read using
pairwise sequence alignment --- the same kind of comparison familiar from
BLAST output, where two sequences are scored by their percent identity over
their alignable region. There is one critical caveat for readers used to
thinking about ITS identity thresholds at the OTU or species level: those
thresholds (commonly cited at 99.0--99.5\% for ``same OTU'' and around 95\%
for ``related but distinct'') are calibrated for \emph{consensus} sequences,
where nanopore basecalling error has already been averaged out. The
comparisons here are between \emph{raw reads}, each carrying its own
basecaller error, so two reads from the same biological sequence routinely
score in the 90s rather than at 99\%+. The graph is built in raw-read identity
space; the OTU thresholds are not the right yardstick for it.

In a specimen with a few hundred reads this exhaustive comparison is cheap; in
a specimen with tens of thousands of reads it is the dominant cost in the
pipeline, and Section~\ref{sec:scalability} describes how 0.7 made it
tractable at scale. The output is a similarity matrix: read $i$ vs read $j$
has score $s_{ij}$.

The matrix is then sparsified into a graph. For each read, only the edges to
its few most-similar neighbors are retained --- this is the
``K-nearest-neighbors'' step, with $K$ controlled by
\texttt{--k-nearest-neighbors} (default 5). Keeping only the strongest edges
does two useful things at once: it removes noise edges between weakly-related
reads that would otherwise pull MCL toward over-merging, and it prevents the
graph from becoming so dense that MCL's matrix operations slow to a crawl.
The remaining edges, weighted by their identity scores, form the graph that
MCL clusters.

The reason this graph structure works for variant discovery is mechanical:
real biological variants in a specimen produce dense subgraphs (each read is
highly similar to other reads from the same variant, and a few hops away from
reads of a different variant), while sequencing noise produces a thin scatter
of edges with no internal structure. MCL is good at finding the dense
subgraphs and ignoring the scatter. The initial proof of concept was
strikingly good at this from the start --- minor sequence variants surfaced
as separate clean clusters even though MCL was only ever shown a sparse K-NN
slice of the pairwise identity matrix. That early result was what convinced
us the rest of the pipeline was worth building.

\section{Variant phasing through MSA}

Speconsense has always done a kind of variant phasing implicitly --- whenever
MCL splits a specimen's reads into multiple clusters, the resulting per-cluster
consensuses already represent distinct variants. What was missing through the
0.4 and 0.5 series was \emph{explicit} variant phasing: a mechanism to discover
variation \emph{within} a cluster and split it on a per-position basis.
Position-based, MSA-driven phasing was prototyped in November 2025
(\texttt{8f18e8d} added variant position detection on November 9;
\texttt{44ab13b} implemented automatic phasing on November 11) and shipped in
v0.6.0 on December 5, 2025.

The approach is straightforward in outline: build a multiple sequence alignment
(MSA) of the reads in a cluster using SPOA~\cite{vaser2017,lee2003}, find columns
where multiple bases each occur with non-trivial frequency, and split the reads
by their pattern across those positions. The version that shipped in v0.6.0
uses a recursive hierarchical algorithm that picks the most informative
position at each level and partitions reads, recursing into the resulting
subsets until no further phaseable positions remain.

For positions that survive variant analysis but cannot cleanly partition the
reads --- typically because the minor allele appears at intermediate frequency
without grouping with other minor alleles --- the consensus emits an IUPAC
ambiguity code (R for A/G, Y for C/T, etc.). Ambiguity calling was added in
\texttt{306a00e} (December 1, 2025) and gave validators a way to see residual
heterogeneity that the phaser couldn't resolve.

\subsection{A note on terminology}

We use ``variant phasing'' rather than ``haplotype phasing'' throughout. Many
of the variants the pipeline finds in fungal ITS data are not haplotypes in the
strict population-genetic sense: paralogous rDNA copies, intragenomic variation
within a single individual, and other non-allelic differences appear regularly,
and ``variant'' covers all of these without overcommitting to a model.

\subsection{Two cluster-level metrics: \texttt{err\_factor} and \texttt{cer\_factor}}
\label{sec:metrics}

Two derived metrics recur throughout the rest of the paper. We introduce them
informally here; precise definitions live in the companion
papers~\cite{cer_paper,cer_in_practice}.

\begin{itemize}
  \item \texttt{err\_factor} is a per-cluster quality metric. It is the ratio
  of the disagreement actually observed among the reads of a cluster to the
  disagreement an empirical basecaller error model predicts for that cluster.
  A cluster of reads that are all faithful copies of a single biological
  sequence, differing only by basecaller noise, has \texttt{err\_factor}
  $\approx 1$. A cluster that is a mix of related-but-distinct sequences has
  \texttt{err\_factor} much greater than 1, because the mix produces
  per-position disagreement above what the error model alone would explain.
  \texttt{err\_factor} is \emph{peer-independent}: it can be computed from a
  single cluster in isolation.

  \item \texttt{cer\_factor} is a peer-relative significance metric drawn
  from the \emph{critical error rate} (CER) framework~\cite{cer_paper,cer_in_practice}:
  a per-variant significance test calibrated against the platform's
  actual error rate, named for the per-position rate
  $p^\ast$ at which an observed variant pattern would just become
  plausible as artifact. For a smaller cluster sitting near a larger
  peer in the same identity group, \texttt{cer\_factor} asks: what
  per-position error rate would the basecaller need for the smaller
  cluster's variant pattern to be plausibly explained as basecaller
  artifact, expressed as a multiple of the basecaller's actual error
  rate? \texttt{cer\_factor} much less than 1 means the actual basecaller
  rate is already higher than what would be needed --- the variant
  pattern is well within what basecaller noise could produce, and is
  most likely noise. \texttt{cer\_factor} much greater than 1 means the
  basecaller would need to err many times faster than it actually does
  to produce the observed pattern --- the variant is implausible as
  basecaller artifact and is likely real biology. The dimensionless
  factor is $p^\ast/q_\text{ctx}$ in the notation of~\cite{cer_in_practice}.
  We will refer to ``CER'' and ``the significance framework''
  interchangeably from here on.
\end{itemize}

The two metrics answer two genuinely different validator questions: ``is this
cluster real?'' (\texttt{err\_factor}) and ``is this variant distinguishable
from its peer?'' (\texttt{cer\_factor}). The 0.8 pipeline computes and surfaces
both, and uses each for the decision it is appropriate for.

\section{From proof of concept to production: 0.4 through 0.6}

Through the spring 2025 dormancy, the tool was a proof of concept ---
runnable but not used by anyone other than the author. Two events in the
autumn brought it into active partnership. Following Harte Singer's
August enquiry, Stephen Russell at Mycota Lab began testing speconsense
seriously against production data in early October 2025. A detailed UX conversation in mid-October turned into a release one week
later: \textbf{v0.4.0 (October 23, 2025)} introduced customizable FASTA header
fields, a quality report aimed at flagging sequences worth a second look,
MSA-based variant merging with \texttt{.raw} traceability files for full
provenance, and incremental output writing for large datasets. The release
addressed most of the items on Stephen's feedback list and marked the
transition from proof-of-concept to \textbf{preview} --- a stable enough
surface for someone other than the author to run real workloads on it.

Through late October and November, Mycota Lab ran experimental batches with
v0.4.x and v0.5.x. Two pieces of feedback from this period drove the next
major release.

The first (October 30) was a feature request for primer-pool support. Mycota
Lab's workflow amplifies a single specimen with multiple primer pairs that
target overlapping but non-identical stretches of the same locus --- in this
case ITS1F/ITS4 (full ITS) and gITS7/ITS4 (ITS2 only, with the forward primer
sitting in the 5.8S region). The two amplicons share the ITS2 + LSU end but
differ in length, and the summarize tool needed to recognize when a $\sim$360
bp gITS7--ITS4 amplicon and a $\sim$660 bp ITS1F--ITS4 amplicon were ``the
same sequence'' with respect to their shared region. Implementation was
deferred to a fast-follow release, \textbf{v0.6.1} (Section~\ref{sec:chimeric}).

The second (November 7) was a flag from Stephen that residual unphased
heterogeneity was visible in some specimens --- clear two-base patterns at
certain positions across most of a cluster, but the consensus reporting only
the majority. This aligned with MSA analysis work the author had begun the
same week, and turned into the headline feature of the next major release.
\textbf{v0.6.0 (December 5, 2025)} shipped automatic detection of variant
positions, MSA-based variant phasing, and IUPAC ambiguity codes for positions
that didn't cleanly resolve into separate variants. v0.6.0 was the first
release used in production at Mycota Lab.

The principle behind the ambiguity-calling work was \textbf{always express
uncertainty clearly}. If the reads disagree at a position and the disagreement
cannot be cleanly resolved by phasing, the consensus should say so --- by
emitting an IUPAC ambiguity code (R for A/G, Y for C/T, etc.) rather than
picking a winner. The previous behavior had let the consensus assert certainty
(a single base) where the underlying reads showed none, which validators
reasonably read as the tool actively hiding information.

This was the right principle but the implementation was under-defended.
Ambiguity calls were generated at any variant position that survived the
frequency threshold but failed to phase, with no separate test for whether the
cluster as a whole was clean enough for an ambiguity call to be meaningful.
Validators eventually started seeing the symmetric problem: clusters with so
many ambiguity codes that the consensus was unhelpful, because the cluster
itself was a mix of related-but-distinct sequences that should have been
split apart or filtered out. The principle was correct, but it needed
noise-reduction machinery that did not arrive until 0.8.0
(Section~\ref{sec:hp-and-ambiguity}). In the meantime, ambiguity codes did at
least make the underlying issue visible, which is more than the previous
behavior had done.

\subsection{The chimeric merge issue (v0.6.1)}
\label{sec:chimeric}

The primer-pool feature requested in October shipped in v0.6.1 (December 18,
2025) as overlap-aware identity scoring: the distance metric between two
sequences would only score the region they had in common. This worked for the
intended primer-pool case. It also created a chimeric-merge failure mode:
two sequences from genuinely different organisms that happened to share enough
overlap with a third sequence could merge transitively into a chimera. The
fix, in \texttt{3f82c66}, made the overlap distance fire only between
sequences with \emph{different} primer pairs, restoring strict identity
scoring within a primer pair. The current rule is still permissive: any two
amplicons from different primer pairs are eligible to merge on their
overlapping region. Future work will tighten this to merge only across
primer pairs whose loci are \emph{a priori} expected to overlap biologically
(e.g., an ITS1F/ITS4 amplicon and a gITS7/ITS4 amplicon), refusing the merge
when the primer-pool relationship does not predict biological co-occurrence
of the shared region.

\subsection{IUPAC conflation bugs}

Two IUPAC bugs surfaced in the 0.6 series, with different origins but the same
flavor.

The first (\texttt{ca3fc20}, v0.6.4) was a code-expansion bug in variant
merging. When merging a base with an existing IUPAC code --- say C against Y
--- the code naively unioned the literal symbols and produced N, when the
correct result was Y (since Y already covers C and T). The fix added a helper
that expanded existing codes to their constituent bases before re-encoding.

The second (\texttt{b83449f}, December 2025 / v0.7) was a defaults issue, not
a code bug, but presented to users the same way. Summarize would merge
variants and propagate IUPAC codes regardless of the relative sizes of the
input clusters. A 3-read cluster with one minor SNP could merge into a
500-read cluster and inject an ambiguity code into an otherwise clean
consensus. The validators' implicit expectation was that every \emph{read}
would count equally; the implementation gave every \emph{cluster
consensus} an equal vote, which is a different thing. The design
rationale was not equal voting per se, but \emph{preservation of
information} --- when a small cluster carried a SNP the dominant cluster
did not, we wanted that difference visible in the merged consensus
rather than silently dropped. The two intentions agree most of the time
but diverge when cluster sizes are imbalanced, which is exactly when an
ambiguity code in an otherwise-clean consensus surprises the reader. The
fix raised the default \texttt{--merge-min-size-ratio} from 0.0 to 0.1,
requiring a candidate to be at least 10\% the size of the dominant
cluster to participate in the merge.

\section{Scalability and ergonomics: 0.7}
\label{sec:scalability}

The 0.7 series (v0.7.0 in January 2026 through v0.7.7 in March 2026) was the
operational hardening release. The pieces:

\begin{itemize}
  \item \textbf{Scalability mode.} Many of the pairwise operations in the
  pipeline are O($n^2$) by default, which is fine at hundreds of reads and
  painful at thousands. A vsearch-backed~\cite{vsearch} candidate finder
  reduces the dominant operations to roughly O($n \log n$) by pruning the
  candidate set before edlib~\cite{edlib} scoring. Activates automatically
  above a configurable threshold.
  \item \textbf{Profile system.} YAML presets (\texttt{herbarium},
  \texttt{strict}, \texttt{nostalgia}, \texttt{largedata}, \texttt{compressed})
  expose a small set of curated parameter combinations matched to common
  workflows, with user-editable profiles in
  \texttt{\textasciitilde/.config/speconsense/profiles/}.
  \item \textbf{Threading.} A \texttt{--threads} option parallelizes consensus
  generation and clustering across cores.
  \item \textbf{Subpackage refactor.} The flat \texttt{core.py} and
  \texttt{summarize.py} were split into the current \texttt{core/} and
  \texttt{summarize/} subpackages, both for readability and to make selective
  testing tractable.
\end{itemize}

By the end of the 0.7 series the pipeline was stable and the main remaining
concerns came from the validator community.

\subsection{The 0.7 pipeline at a glance}

For reference when reading the 0.8 description in
Section~\ref{sec:pipeline-08}, the 0.7 pipeline ran six phases per specimen:

\begin{center}
\begin{tabular}{cl p{0.55\linewidth}}
\toprule
Phase & Name & What it does \\
\midrule
1 & Initial clustering & MCL on a K-NN identity graph (greedy fallback) \\
2 & Pre-phasing merge & Combine clusters with HP-equivalent consensuses \\
3 & Variant phasing & MSA-based variant splitting; emit IUPAC codes for unphased positions \\
4 & Post-phasing merge & Combine subclusters that became HP-equivalent after phasing \\
5 & Size filtering & Drop clusters below \texttt{--min-size} and \texttt{--min-cluster-ratio} \\
6 & Output generation & SPOA consensus, primer trimming, FASTA/FASTQ writeout \\
\bottomrule
\end{tabular}
\end{center}

The summarize tool then ran identity grouping, MSA-based variant merging,
optional cross-primer overlap merging, variant selection, and a quality report.
Two architectural points to retain for the 0.8 comparison: (a) the levers
available for variant selection were ad hoc --- absolute size floor, ratio to
the dominant cluster --- with no statistical grounding, so a 5-read variant in
a 300-read specimen and a 5-read variant in a 30-read specimen were treated
identically despite the very different per-cluster significance; and (b)
summarize made its own grouping decisions independent of core's clustering.

\section{The three problems behind 0.8.0}
\label{sec:hp-and-ambiguity}

Three classes of validator question kept recurring through the 0.7 series.
Each one pointed at an architectural decision rather than a localized bug.

\subsection{What to do about homopolymers}

The most frequent question was a variant of: ``I have a 50/50 mix of reads
that look like CCCTT and CCTTT. The middle position looks like a clean SNP ---
why doesn't speconsense call an ambiguity here?'' This was raised
independently and over many specimens by Stephen Russell (Mycota Lab) and by
Elora Adamson, Local Coordinator for the MycoMap BC Network project (British
Columbia); Elora provided the most detailed cluster-level feedback and a
running set of specimens that exhibited the issue.

The interface-homopolymer pattern was the recurring case study, but the
underlying question is broader: what should the pipeline do about homopolymers
\emph{at all}? Through 0.7, speconsense applied homopolymer (HP)
normalization whenever it compared two cluster \emph{consensus} sequences
--- in pre- and post-phasing merges, identity-group formation, and
summarize's variant merging. (Read-vs-read comparisons during initial
similarity-matrix construction use raw edlib identity and have never been
HP-normalized.) The mechanism in core is a native HP-aware scoring loop
on top of an edlib alignment: within an HP run of qualifying length, an
insertion or deletion of the run's base is not penalized, so a length-7
A-run and a length-8 A-run align as equivalent. The summarize tool uses
the related \texttt{adjusted-identity} package~\cite{adjusted_identity}
for the same purpose with a slightly different parameterization. In
either implementation, runs of identical bases are not literally
compressed to single characters, but the effective scoring is as if they
were. Under that lens, CCCTT and CCTTT are equivalent --- the only
difference between them is the length of the C run and the length of
the T run, and homopolymer length differences are the dominant systematic
error mode in nanopore sequencing. The pipeline was designed to ignore them.

But the SNP interpretation is also reasonable. The middle position genuinely
is C in some reads and T in others. A validator looking at a stack of reads
doesn't see ``two homopolymer runs of variable length'' --- they see a column
where two bases appear at meaningful frequency, and ask why their tool isn't
flagging it.

The right answer is empirical: how often is the SNP interpretation correct,
and how often is it the homopolymer-noise interpretation?
Paper~2~\cite{hp_paper} addresses this directly. Briefly: the rate of
artifactual variation at homopolymer-interface positions is no higher than at
non-homopolymer positions (rate ratios of 1.03 and 0.88 for Dorado v3.5 and
v5.0 respectively), so blanket normalization at short homopolymer lengths is
too aggressive --- it hides real variation behind a noise model that doesn't
actually apply at those lengths. The fix is not to abandon HP normalization
(it remains correct at long HP lengths, where the error rate genuinely is
higher) but to make it context-aware: a per-event-type, per-HP-length error
rate, applied as a continuous statistical weight rather than a hard switch.

Relaxing the normalization is mechanically straightforward but required new
infrastructure: a context-aware error model (per-basecaller $q_\text{ctx}$
tables~\cite{hp_paper}, keyed by event type and homopolymer run length) so
that variants at HP positions could be evaluated statistically rather than
reflexively suppressed. We will refer to this empirical error model simply
as the \emph{error model} from this point on. The same model also drives
the cluster-homogeneity metric \texttt{err\_factor} (Phase~12) and the
peer-relative \texttt{cer\_factor} (Phase~10), so the work that started as
``let validators see HP variation when it's real'' ended up enabling the
cluster-quality machinery as well.

\subsection{Ambiguity inflation in low-quality clusters}
\label{sec:ambiguity-inflation}

The second recurring question was the inverse: ``Why does this consensus have
so many ambiguity codes? It looks like noise.''

Investigation showed that some clusters were genuinely heterogeneous --- a mix
of reads from related but distinct sequences that the clustering step had
grouped together because they were too close to separate at the configured
identity threshold. Variant phasing would then attempt to split them, often
fail to find a clean partition, and emit ambiguity codes at every disagreeing
position. A 600 bp consensus with 30 ambiguity codes is unhelpful to a
validator: it's neither a clean sequence nor a clear signal that the cluster
should be discarded. There was no way for the tool to say ``this cluster is
not a real sequence, it is a noise blob'' --- the only filter was on size,
which is fine for the easy cases but is a poor proxy when the goal also
includes rescuing low-read specimens. A noise blob that survives the size
floor is hard to flag without a distinct test for cluster cohesion.

\subsection{Validating low-support variants}

The third question was the hardest: ``I have a 5-read cluster that's two SNPs
different from a 200-read cluster. Is it a real variant?''

Through 0.7, the answer was effectively ``we don't know --- that's your call.''
Variants were filtered by absolute size and by ratio to the dominant cluster,
and variants that survived both were emitted, but neither lever is
statistically optimal across a dataset with a 100$\times$ or wider range in
per-specimen read depth. A 5-read variant in a low-error context might be much
more credible than a 50-read variant in a length-7 G homopolymer, but the tool
had no way to express the difference.

This third concern motivated the variant significance framework
(Paper~1~\cite{cer_paper}) and its empirical extension
(Paper~3~\cite{cer_in_practice}). The mathematical content is in those
papers; the pipeline machinery needed to act on the framework --- a
chain of upstream cleanup phases plus the CER test itself --- is the
subject of Section~\ref{sec:pipeline-08}. The caveat below applies to
the framework as a whole.

A careful caveat is in order before proceeding. Even with the CER framework,
the pipeline cannot say which low-support variants are real --- it can only
say which ones are not artifacts of nanopore basecaller error. A variant
that the framework rejects (\texttt{cer\_factor} below threshold) is not
necessarily wrong biology; it might be real biology that simply lacks the
read support to distinguish from basecaller noise on this run. And a variant
that passes the threshold is not necessarily real biology either; it might
be a chimera, a contaminant, an upstream demultiplexing error, or any
artifact whose generative process is something other than basecaller error.
What the framework gives us is a clean separation along a single axis:
``could this be basecaller noise?'' --- yes or no. Everything else remains
the validator's call. In practice, validators answer those other
questions through long-term, multi-run accrual of reference data:
variants that recur across specimens and across sequencing runs
accumulate evidence for biological reality that no single specimen's
read pile can supply. The CER framework is a per-run filter; the
validator's reference-database work is the cross-run integrator.

\section{The 0.8.0 pipeline}
\label{sec:pipeline-08}

The 0.8.0 pipeline runs in thirteen sequential phases. Six are new in
0.8.0: Phases 5--10, a coupled cleanup-and-routing cascade. Phase~9's
MAD-based outlier removal shares its intention with 0.7's
static-threshold outlier removal --- pull a small tail of misfit reads
out of an otherwise-clean cluster --- but the underlying mechanism is
different enough that it behaves like a new phase, and
§\ref{sec:results} shows it is the single largest quality contributor
in the additive ladder. The remaining phases are largely unchanged
from 0.7 in mechanism but renumbered into the flat sequence. This
section describes each phase and the design intent behind it;
quantitative results live in Section~\ref{sec:results}.

The new phases form a connected cascade rather than a list of independent
improvements. The CER framework (Phase~10) was the original target --- a
defensible statistical test for low-support variants. Implementing it required
confidence that the variant positions being tested were real, which surfaced a
chain of upstream issues. We needed the consensus that drives variant-position
analysis to be trustworthy (Phase~5: noise filter --- disband clusters whose
consensus has positions with no majority). We needed reads to be assigned to
the cluster they best fit, not merely the one MCL placed them in (Phase~6:
read reassignment), and we needed reads earlier phases had discarded to get a
second chance once the cluster landscape stabilized (Phase~7: discard
recovery). When discard recovery brought new reads into existing clusters, it
exposed variant positions the first phasing pass had missed (Phase~8: second
phasing pass). Even after all of that, a small number of outlier reads in
otherwise-clean clusters were inflating residual ambiguity, motivating a
per-cluster cleanup (Phase~9: MAD-cleaned consensus generation) that
runs immediately before CER so the statistical test sees the
stabilized cluster state.
The phases were designed together, and we will see in
Section~\ref{sec:results} that they only deliver their intended
benefit \emph{together}: turning some on without their upstream
prerequisites can actively hurt.

\subsection*{Phase 1: Initial clustering}

MCL on a K-NN similarity graph by default; greedy fallback if MCL is
unavailable. Unchanged in mechanism since 0.6. Immediately after Phase 1,
a \emph{hard floor} of 3 reads is applied: clusters with fewer than 3 reads
are dropped to discards, since reliable error correction requires at least
three observations per position. This hard floor is not a configurable knob
and applies regardless of the user-set \texttt{--min-size} (which can be set
as low as 3 but no lower). Its purpose is structural rather than
quality-driven: a 1- or 2-read ``cluster'' is not a consensus in any
meaningful sense.

\subsection*{Phase 2: Pre-phasing merge}

Combine clusters whose consensuses are equivalent under homopolymer
normalization, applied via core's native HP-aware comparator with a
configurable HP-length threshold (\texttt{--hp-normalization-length},
default 6). HP runs of length $\geq 6$ are treated as HP context (their
length variation is suppressed); shorter runs score normally, so a
length-3-vs-length-4 difference is no longer absorbed silently. This
catches cases where the initial clustering placed reads into separate
clusters that resolve to the same biological sequence after consensus.

\subsection*{Phase 3: Variant phasing}

For each cluster, build an MSA, find variant positions where qualifying bases
reach the configured frequency threshold, and split reads by their allele
combination. A cluster with no variant positions passes through unchanged.
Three changes in 0.8.0 sharpen this step relative to 0.7: (a) a position
can now qualify either across the full cluster \emph{or} within a
quality-biased sample (top-quality reads), which catches signal that
low-quality reads were diluting out; (b) the SPOA alignment parameters are
tuned to reduce the relative gap penalties at this step, which makes it
easier for the alignment to distinguish a true substitution from a
homopolymer-length artifact and exposes more variant positions cleanly; and
(c) reads that do not match either the major or a phased minor allele
cleanly are now sent to the discard pool for later recovery (Phase~7),
where 0.7 had instead force-assigned them to the closest-matching
sub-cluster (commit \texttt{372fa5b}, April 2026). The previous behavior
quietly contaminated otherwise-clean sub-clusters with reads that matched
neither allele well and surfaced as ambiguity codes in the consensus.

\subsection*{Phase 4: Post-phasing merge}

Run the pre-phasing merge logic again --- same HP-length threshold, same
HP-aware comparator --- this time on the subclusters produced by phasing.
Phasing can produce clusters that are HP-equivalent to clusters elsewhere
in the specimen; this consolidates them.

\subsection*{Phase 5 (NEW): Noise filter}

Run SPOA on each remaining small cluster. If the resulting MSA has columns
where no base accounts for more than half of the reads, the cluster is
disbanded --- its consensus is not really a consensus. Concretely, a 4-read
cluster with a 2-2 split at some position fails the test, as does a 3-read
cluster with three different bases at a position; both patterns are common
in genuinely noisy small clusters. Reads from disbanded clusters fall to the
discard pool and may return via Phase~7.

\emph{Why this step}: A cluster whose own consensus is unreliable cannot serve
as a reference for downstream variant-position analysis (Phases~6, 7, 8) or
pairwise statistical evaluation (Phase~10); better to disband it and let its
reads find a more cohesive home in Phase~7. A second mechanism, less
obvious but empirically substantial: noisy small clusters tend to generate
spurious variant positions during MSA analysis, which then feed Phase~6's
concordance scoring as if they were real signal. Removing those clusters
prunes the false-positive variant positions before reassignment uses them.

\subsection*{Phase 6 (NEW): Read reassignment}

Within each complete-linkage identity group, test whether each read agrees
better with its current cluster's consensus or with another cluster's
consensus in the same group. If concordance with a different cluster is
meaningfully better, move the read. An edit-distance guard prevents reads
from drifting across distinct sequences. Reassignment uses both edit distance
and per-position concordance at variant sites --- the same machinery that
Phase~7 will reuse.

\emph{Why this step}: MCL clustering is good but not perfect; reads can land
in a sibling cluster they don't quite fit. The signal that lets us catch
this in Phase~6 was not available in Phase~1: at initial clustering time we
do not yet know where the meaningful variant positions are, so we have no
better basis for read-to-cluster comparison than raw edit distance. By
Phase~6, variant phasing has surfaced concrete variant positions, and we
can compare each read's bases at those positions against each candidate
cluster's consensus --- a much more discriminating signal. Misassigned
reads contribute their disagreements to the wrong consensus, inflating
ambiguity; letting reads migrate to the cluster that best explains them at
its variant positions tightens each cluster's internal consistency and
reduces ambiguity counts.

\subsection*{Phase 7 (NEW): Discard recovery}

Reads that were dropped earlier --- failed phasing, MAD outlier removal,
hard-floor filters, Phase~5 noise filtering --- are tested against the
surviving cluster consensuses. Reads that are concordant with a surviving
consensus rejoin it. Auto-skipped if Phase~6 is disabled, since the two
phases share machinery.

\emph{Why this step}: Reads discarded earlier weren't necessarily bad reads
--- they often just didn't fit any cluster well at the time, because the
cluster landscape was still in flux. After Phases~5 and~6 stabilize the
surviving clusters, many discarded reads do fit a surviving consensus
cleanly; admitting them increases support for real variants without changing
what variants exist.

\subsection*{Phase 8 (NEW): Second phasing pass}

If Phases~6 or~7 changed any cluster's membership, that cluster's MSA
structure may have shifted enough to support a phasing decision that was not
previously available. Re-run Phase~3 on those clusters. Gated by both
\texttt{--enable-position-phasing} and \texttt{--enable-read-reassignment}
--- without either, nothing changed and there is no work to do.

\emph{Why this step}: Discard recovery (Phase~7) routinely brings reads into
existing clusters that surface variant positions the first phasing pass
missed --- either because the alternate allele was below the count threshold,
or because the supporting reads weren't yet present in the cluster. Without a
second pass, this newly-visible variation collapses to ambiguity codes in the
consensus instead of resolving into separate variants. Phase~8 is primarily
a quality tool (ambiguity reduction) and secondarily a discrimination tool
(variant discovery): the same variation that fails to phase the first time
through, when surfaced cleanly on the second pass, both reduces residual
ambiguity in the larger cluster and produces a separate variant for the
validator's review.

\subsection*{Phase 9 (NEW): MAD-cleaned consensus generation}

Build an initial consensus for each surviving cluster, then remove reads
whose identity to that consensus is more than $n \times$ MAD from the
median --- a small number of outlier reads can otherwise inflate
the per-cluster disagreement reading on a cluster that is biologically
clean. Rebuild the consensus on the surviving reads. \emph{MAD} stands
for median absolute deviation --- for a cluster of read-vs-consensus
identities, the MAD is the median of the absolute differences between
each read's identity and the cluster's median identity. Reads whose
identity sits more than several MADs away from the median are
statistical outliers under any reasonable assumption about the
underlying distribution, regardless of whether the cluster's reads
are 99\%-identity-clean or 92\%-identity-noisy; this scale-free
property is why the previous static \texttt{--outlier-identity}
threshold is gone. Runs unconditionally; four CLI parameters
control which reads are flagged (modified-Z threshold, gap rule, MAD
floor, safety floor on absolute drop). §\ref{sec:looking-ahead}
reports the 0.8.1 tuning of these parameters.

\emph{Why this step}: Even after Phases~5--8 stabilize the cluster
landscape, a small number of outlier reads in otherwise-clean clusters
keep injecting per-position disagreement that downstream phases would
otherwise see as residual heterogeneity. Running this cleanup before
CER (Phase~10) and \texttt{err\_factor} computation (Phase~12) means
both downstream tests see a settled cluster state: CER's $M$ and $N$
read counts are post-MAD, and \texttt{err\_factor} measures
disagreement on the cleaned cluster. Without MAD, \emph{any} choice
of \texttt{--max-err-factor} is less discriminating: the threshold
cannot cleanly separate completely problematic clusters from clusters
that are good but happen to carry a few outlier reads, because both
populations crowd into the same \texttt{err\_factor} range. MAD pulls
them apart, restoring the threshold's resolving power.

\subsection*{Phase 10 (NEW): CER validation}

Within each identity group, the largest cluster is designated the
\emph{anchor}: it is presumed to represent the dominant biological signal in
the group, and is exempt from significance testing. For every non-anchor
cluster in the group, find its larger peers (the anchor and any other
larger non-anchor) and compute the per-position context-aware critical
error rate (CER) factor for each pair. Annotate the cluster with the
minimum factor across peers --- its weakest case. Anchors and clusters with
no comparable peer carry \texttt{cer\_factor=None} and are exempt.
Crucially, this phase only \emph{annotates}. The pass/fail decision is
applied later, by \texttt{speconsense-summarize}.

\emph{Why this step}: The original target of the 0.8 redesign. The other new
phases inflate the variant count (more variants surfaced from cleaner
clusters), which would in turn inflate validator review burden if every
variant passed through to the main output. CER provides the per-variant
statistical test that lets the pipeline say ``this small variant is almost
certainly a basecaller-noise replica of a larger peer,'' routing
demonstrably-bad variants to a separate track instead of dropping them.

\subsection*{Phase 11: Size filtering}

Drop clusters below \texttt{--min-size} or \texttt{--min-cluster-ratio}.
The Phase~1 hard floor at 3 reads still applies independently as a lower
bound. The mechanism is unchanged from 0.7, but the defaults are not: 0.7.7
shipped \texttt{min-size=5} and \texttt{min-cluster-ratio=0.01}, and 0.8.0
relaxes both to \texttt{min-size=3} and \texttt{min-cluster-ratio=0}. The
relaxation is intentional: the discrimination work that previously had to
happen at this size-floor step has moved upstream to the noise filter
(Phase~5) and downstream to CER routing (Phase~10) and the
\texttt{err\_factor} filter (Phase~12). Phase~11 in 0.8 is a safety net,
not a primary filter.

\subsection*{Phase 12 (NEW): Output emission}

Write the final FASTA, FASTQ, and MSA for each surviving cluster. Compute
\texttt{err\_factor} (the cluster's observed-vs-expected disagreement
ratio under the basecaller error model) on the post-MAD MSA from
Phase~9, and stamp it into the FASTA header for downstream routing in
\texttt{speconsense-summarize}.

\emph{Why this step}: Routing decisions (which records pass, which route
to \texttt{.lq}) need a peer-independent quality metric, and
\texttt{err\_factor} provides it. Because the MSA was already built
during the post-MAD consensus pass, no additional alignment work is
needed here --- emission is a thin wrapper over an existing artifact.

\subsection*{Phase 13: Discard reads written}

When \texttt{--collect-discards} is set, write the discard pool to disk for
inspection. Mechanism unchanged from 0.7.

\subsection{\texttt{err\_factor} and the \texttt{.ns}/\texttt{.lq} routing}

Returning to the two cluster-level metrics introduced in
Section~\ref{sec:metrics}, the 0.8 summarize tool routes records on each:

\begin{itemize}
  \item In \texttt{speconsense-summarize}, records below
  \texttt{--min-cer-factor} (default 1.0) route to \\
  \path{__Summary__/variants/{name}.ns-RiC{ric}.fasta}.
  \item Records above \texttt{--max-err-factor} (default 1.5) route to \\
  \path{__Summary__/variants/{name}.lq-RiC{ric}.fasta}.
  \item \texttt{.lq} takes precedence when both fire. Neither filter drops
  records --- they redirect them to a separate inspection track for
  later review.\footnote{In current practice, the lab sequencer running
  the primary data analysis (e.g., Stephen Russell at Mycota Lab, Harte
  Singer at FUNDIS) is the immediate consumer of the \texttt{.ns} and
  \texttt{.lq} tracks; downstream validators work on MycoMap, which at
  the time of writing exposes only the pass-track summary and not the
  \texttt{variants/} subdirectory. Surfacing the routed tracks through
  MycoMap is on the roadmap.}
\end{itemize}

In the empirical results below, the two tracks behave as different
inspection surfaces, not as parallel biological-recovery channels.
\texttt{.ns} is dominantly basecaller-noise replicas of pass-track
variants (§\ref{sec:results-routing}: median \texttt{cer\_factor} =
0.022). \texttt{.lq} is empirically a \emph{within-specimen}
variant-exploration track: in our 955-specimen run it adds zero
specimens to pass-track recall (§\ref{sec:results-tax}). Validators
who look only at the pass track miss the within-specimen variant
exploration but see no specimen-level recall loss; validators who
look at \texttt{.lq} see additional variants for already-recovered
specimens but no new organisms.

\section{Empirical results}
\label{sec:results}

We report results from two empirical matrices, both run against the full
ont98 ITS dataset (955 demultiplexed specimens; full methods in
Section~\ref{sec:methods}).\footnote{The ablation matrices in
§\ref{sec:results-additive}--\ref{sec:results-permissivity} and the
cross-tool tables in §\ref{sec:results-ngsi}--\ref{sec:results-tax}
were generated under the pre-reorder pipeline, where CER (then
Phase~9) annotated clusters before MAD outlier removal (then
Phase~11). On the post-reorder \texttt{v08-defaults} re-run reported
in §\ref{sec:results-version}, the headline pass-record / ambiguity /
yield numbers shift by 2--3\% (and the .lq count by under 1\%),
without altering specimen-level recall. Full matrix re-derivation
under the post-reorder ordering and the 0.8.1 MAD defaults is
deferred to a future revision pass; the qualitative findings here
hold under both.} The first is an \emph{additive ladder} on the
v0.8.0 binary: starting from a stripped-down baseline (Phases 1, 2, 3, 4,
11, and 12, with Phase 9 MAD disabled), the new phases are turned on one at a time in
dependency order. Each step measures the marginal contribution of one
phase given everything that came before. The second is a
\emph{cross-version} comparison: v0.7.7 with its own defaults vs.~v0.8.0
with its own defaults, the picture a validator sees when they upgrade.

\subsection{Per-step contributions of the new phases}
\label{sec:results-additive}

The additive ladder, on the validator-facing pass track. Baseline
\texttt{B} = Phases 1, 2, 3, 4, 11, and 12-without-MAD enabled (the
v0.7-equivalent core); each subsequent row turns on one new phase in
dependency order, with \texttt{+MAD} (=v0.8) the terminal step. Each
delta is the marginal contribution of one phase \emph{given everything
before it}; phases interact, so a different ordering would produce
different intermediate values --- §\ref{sec:results-cascade}
shows the alt-order trajectory and confirms that adding P7 before P5
actively hurts before P5 catches up.

\begin{center}
\begin{tabular}{l rr rr rr r}
\toprule
Step & \multicolumn{2}{c}{records} & \multicolumn{2}{c}{yield (reads)} & \multicolumn{2}{c}{ambig codes} & ambig \\
& count & $\Delta$ & count & $\Delta$ & count & $\Delta$ & /rec \\
\midrule
\texttt{B} (baseline) & 3,767 & --- & 315,128 & --- & 1,108 & --- & 0.294 \\
\texttt{+P6} & 3,552 & $-5.7\%$ & 317,451 & $+0.7\%$ & 1,043 & $-5.9\%$ & 0.294 \\
\texttt{+P6+P5} & 3,246 & $-8.6\%$ & 314,816 & $-0.8\%$ & 650 & $\mathbf{-37.7\%}$ & \textbf{0.200} \\
\texttt{+P6+P5+P7} & 2,928 & $-9.8\%$ & 371,752 & $\mathbf{+18.1\%}$ & 973 & $+49.7\%$ & 0.332 \\
\texttt{+P6+P5+P7+P8} & 3,355 & $+14.6\%$ & 366,877 & $-1.3\%$ & 799 & $\mathbf{-17.9\%}$ & 0.238 \\
\texttt{+MAD} (=v0.8) & 3,689 & $+10.0\%$ & 366,625 & $-0.1\%$ & 927 & $+16.0\%$ & 0.251 \\
\bottomrule
\end{tabular}
\end{center}

The same ladder, viewed through the above-expected-disagreement lens
(percentage of pass-track records and yield with \texttt{err\_factor}
$>$ 1):

\begin{center}
\begin{tabular}{l rr rr}
\toprule
Step & \multicolumn{2}{c}{records err$>$1 (\%)} & \multicolumn{2}{c}{yield err$>$1 (\%)} \\
& share & $\Delta$ & share & $\Delta$ \\
\midrule
\texttt{B} (baseline) & 24.5 & --- & 6.6 & --- \\
\texttt{+P6} & 24.7 & $+0.8\%$ & 7.0 & $+6.1\%$ \\
\texttt{+P6+P5} & 21.4 & $\mathbf{-13.4\%}$ & 6.5 & $-7.1\%$ \\
\texttt{+P6+P5+P7} & 27.1 & $+26.6\%$ & 9.7 & $\mathbf{+49.2\%}$ \\
\texttt{+P6+P5+P7+P8} & 26.0 & $-4.1\%$ & 8.7 & $-10.3\%$ \\
\texttt{+MAD} (=v0.8) & \textbf{16.6} & $\mathbf{-36.2\%}$ & \textbf{2.6} & $\mathbf{-70.1\%}$ \\
\bottomrule
\end{tabular}
\end{center}

Reading the table: the ambig/rec column is the per-record cleanup
rate, removing the confound that record counts shift between steps.

\begin{itemize}
\item \textbf{P5 (noise filter)} cuts pass-track ambiguity by $-37.7\%$
when added to a pipeline that already has P6: it is directly removing
noise-blob clusters whose ambiguity codes were polluting the pass track.
At the core source level, the same step cuts ambiguity by $-84\%$ (3,820
$\to$ 375 ambig codes per core record); most of P5's benefit happens
upstream of the routing filters.

\item \textbf{P7 (discard recovery)} is the yield engine: $+18.1\%$
pass-track yield, broadly distributed (676 of 828 specimens gain reads,
only 81 lose). It does not significantly change variant count --- the
recovered reads consolidate into existing clusters --- but its raise in
ambig (+49.7\%) is the cost: some recovered reads bring residual
disagreement into otherwise-clean consensuses. The quality lens shows
the same pattern: above-expected-disagreement yield jumps from 6.5\%
of pass to 9.7\% as the recovered reads inflate the noise share.

\item \textbf{P8 (second phasing pass)} resolves much of that residual
disagreement: $-17.9\%$ ambig and $+14.6\%$ records. Variation that was
collapsing to ambiguity codes in larger clusters is now phased into
separate, smaller variants for validator review.

\item \textbf{MAD outlier removal} is the dominant single quality
contributor in the ladder. Its $+10\%$-records bump is not new
discovery: the per-cluster \texttt{err\_factor} distribution tightens
enough that records that previously routed to \texttt{.lq} now make it
to pass. The above-expected-disagreement share drops from $26.0\%$ to
$16.6\%$ of records ($-9.4$pp) and from $8.7\%$ to $2.6\%$ of yield
($-6.1$pp; a $-70\%$ relative reduction) in this single step. The new
phases by themselves leave the noise share roughly flat against
baseline; MAD does most of the quality lifting in the additive
sequence.
\end{itemize}

A sanity check confirms the ladder is consistent with the original
ablation matrix: \texttt{+P6+P5+P7+P8} matches a separate \texttt{v08-no-mad}
ablation run (full 0.8 pipeline minus MAD) to within rounding noise (0
records / +21 reads / +4 ambig codes out of 3,355 records). The same set
of phases enabled produces the same outputs regardless of how they were
arrived at.

\subsection{The new phases are a coupled cascade}
\label{sec:results-cascade}

The most striking result from the additive matrix is what happens when
the phases are added in a different order. We re-ran the ladder with
\texttt{P7} added before \texttt{P5} (the alternate sequence
\texttt{B} $\to$ \texttt{+P6} $\to$ \texttt{+P6+P7} $\to$
\texttt{+P6+P5+P7} $\to \cdots$) and compared the ambiguity-code
trajectory to the linear ordering above:

\begin{center}
\begin{tabular}{l rrrr}
\toprule
ambig codes after $\to$ & \texttt{+P6} & next step & next step & = \texttt{+P6+P5+P7} \\
\midrule
linear order (P5 before P7) & 1,043 & 650 (\texttt{+P5}) & 973 (\texttt{+P7}) & 973 \\
alt order (P7 before P5) & 1,043 & \textbf{1,136} (\texttt{+P7}) & 973 (\texttt{+P5}) & 973 \\
\bottomrule
\end{tabular}
\end{center}

Both paths land at the same end state (the same set of phases enabled
produces the same output, deterministic), but the trajectory is sharply
different. \textbf{Adding P7 to a pipeline without P5 actively increases
ambiguity codes ($+9\%$ above \texttt{+P6} alone).} Discard recovery,
absent the noise-filter prepruning, brings noise-blob reads back into
otherwise-clean clusters where they introduce new ambiguity. Adding P5
afterwards then cleans those up.

The mechanism behind the interaction is concrete: noisy small clusters
generate spurious variant positions during MSA analysis, and those
false-positive positions feed P6's concordance scoring as if they were
real signal. Without P5 to disband the noisy clusters first, P6 reassigns
reads using an inflated set of fake variant positions; P7 then re-admits
discards on the same flawed basis. The combination of P5 + P6 + P7 only
delivers its intended yield-and-ambig benefit when P5 has run first to
prune the false positives.

\textbf{The new phases are not independently optional}: turning on P7
alone is worse than baseline, and the full benefit only appears when
\{P5, P6, P7, P8\} are enabled together as a unit. This is the
empirical fingerprint of the coupled-cascade design we described in
Section~\ref{sec:pipeline-08}.

\subsection{Track separation in the routing dimension}
\label{sec:results-routing}

Phase~10 (CER) and Phase~12 (\texttt{err\_factor}) routing are
summarize-only annotations, independent of the additive core ladder. On
\texttt{v08-defaults} across the 955 ont98 specimens, of 13,300 core
records:

\begin{center}
\begin{tabular}{lrr}
\toprule
track & records & median size (reads) \\
\midrule
pass & 3,689 (28\%) & 8 \\
\texttt{.ns} (CER-rejected) & 7,483 (56\%) & 4 \\
\texttt{.lq} (err-rejected) & 1,659 (12\%) & 3 \\
\bottomrule
\end{tabular}
\end{center}

Every core record is accounted for: 12,831 emit directly to
pass/\texttt{.ns}/\texttt{.lq}, and the remaining 469 are conflated into
414 pass-track records via summarize's identity-group merging. No core
records are dropped. The pass track ends up heavily-supported (mean = 99
reads, p90 = 387 reads); the rejection tracks contain qualitatively
different populations of small clusters.

The \texttt{.ns} track is the strongest single empirical result of the
study. Variants routed there have a \texttt{cer\_factor} distribution
with median \textbf{0.022} and p25 = 0.003: 65\% of \texttt{.ns}-routed
variants have \texttt{cer\_factor} below 0.10 --- orders of magnitude
below the routing threshold of 1.0. The median \texttt{.ns} variant
would need the basecaller's per-position HP error rate inflated 45-fold
to plausibly arise as artifact. The CER framework is not just routing
borderline cases; it is routing variants the math identifies as
essentially-certain basecaller noise.

Phase~9's MAD step shapes the \texttt{err\_factor} distribution and so
shapes the \texttt{.lq} track. Median per-specimen
\texttt{err\_factor\_p90} drops from 2.16 (without MAD) to 1.35 (with
MAD enabled) --- a $-38\%$ reduction in the tail. 76\% of specimens
benefit, 18\% are unchanged, 6\% see the per-specimen p90 inflate (a
metric-shrinkage artifact: MAD removes outlier reads, the smaller
surviving cluster has fewer reads diluting per-position disagreement, so
the disagreement reading concentrates upward on a smaller denominator).
The downstream effects are large: MAD reduces \texttt{.lq} routing by
$-51\%$ (3,390 $\to$ 1,659 records), drops the share of pass records with
above-expected disagreement from $26.0\%$ to $16.6\%$, and drops the
above-expected-disagreement yield share from $8.7\%$ to $2.6\%$ (a
$-70\%$ relative reduction). MAD is, empirically, the largest single
quality contributor among the new phases.

\subsection{Cumulative effect: v0.7.7 \texorpdfstring{$\to$}{->} v0.8.0}
\label{sec:results-version}

Stepping back: what does the cumulative v0.7-to-v0.8 upgrade look like
to a validator running the same data through the new pipeline? Comparing
v0.7.7 with default settings to v0.8.0 with default settings, on the
validator-facing pass track:

\begin{center}
\begin{tabular}{lrrr}
\toprule
metric & v0.7.7 & v0.8.0 & delta \\
\midrule
pass-track variants & 2,831 & 3,609 & \textbf{$+27\%$} \\
pass-track yield (reads) & 348,399 & 366,168 & \textbf{$+5\%$} \\
pass-track ambiguity codes & 1,615 & 902 & \textbf{$-44\%$} \\
records with \texttt{err\_factor} $>$ 1 & 830 (29\%) & 612 (17\%) & \textbf{$-26\%$}\textsuperscript{$\dagger$} \\
yield with \texttt{err\_factor} $>$ 1 & 15,174 (4\%) & 9,402 (3\%) & \textbf{$-38\%$}\textsuperscript{$\dagger$} \\
\texttt{.ns}-routed variants & --- & 7,574 & new \\
\texttt{.lq}-routed variants & --- & 1,665 & new \\
specimens with output & 837 & 850 & $+13$ \\
\bottomrule
\end{tabular}
\end{center}

\noindent\textsuperscript{$\dagger$}\,The \texttt{err\_factor}-related
v0.8.0 rows were derived from the pre-reorder extraction; on the
post-reorder \texttt{v08-defaults} re-run, the pass-record count of
\texttt{err\_factor}~$>$~1 shifts to 558 (a $\sim$0.1pp decrease in
share). A consistent re-derivation for both columns under the
post-reorder ordering and the 0.8.1 MAD defaults is deferred to a
future revision pass.

The headline is one of \emph{clearer separation}: 0.8 emits 27\% more
distinct candidate biological variants per specimen with 44\% fewer ambiguity
codes per consensus, and a small (5\%) yield gain on top. The
\texttt{.ns} and \texttt{.lq} tracks absorb 9,239 additional candidate
variants --- these were either silently averaged away in 0.7 (smearing
the consensus) or filtered out by the per-cluster size threshold. Now
they are visible to validators as separate, statistically-flagged
tracks. Pass-track quality also improves: the count of records with
above-expected disagreement (per-cluster \texttt{err\_factor} $>$ 1,
i.e., observed read-to-consensus disagreement above what the
basecaller error model alone predicts) drops from 830 to 612 records
($-26\%$), and the corresponding yield drops from 15,174 to 9,402
reads ($-38\%$). The $+27\%$ headline is gain in \emph{trustworthy}
records, not just any records.

The headline yield gain of $+5\%$ is heavy-tailed: median per-specimen
delta is $+3$ reads, with 434 specimens gaining (p90 = $+119$ reads) and
376 losing (p10 = $-42$ reads). 0.8 is not uniformly better at recovering
reads; it is better on average and substantially better on a subset.

The headline numbers above blend two distinct kinds of change: the new
phases (mechanism), and a set of relaxed default parameters
(permissivity). Section~\ref{sec:results-permissivity} decomposes the
two; Section~\ref{sec:results-ngsi} sets both against an NGSpeciesID
baseline run on the same dataset.

\subsection{Permissivity vs.~new-phase contribution}
\label{sec:results-permissivity}

The 0.8.0 release ships with three loosened defaults relative to 0.7.7:
\texttt{--min-size} (5 $\to$ 3), \texttt{--min-cluster-ratio} (0.01
$\to$ 0), and \texttt{--hp-normalization-length} (1, all-HP-normalized
$\to$ 6, only HP runs $\geq$ 6 normalized). To isolate the new-phase
effect from the permissivity effect, we run a third configuration,
\texttt{v08-strict}: the v0.8.0 binary with v0.7-style strict thresholds
(\texttt{--min-size 5 --min-cluster-ratio 0.01 --hp-normalization-length
1}). This represents ``the new pipeline judged at 0.7's standard for
what counts as a passing variant.''

\begin{center}
\begin{tabular}{lrrrrr}
\toprule
metric & v0.7.7 & \texttt{v08-strict} & v0.8.0 & new-phase $\Delta$ & permissivity $\Delta$ \\
\midrule
pass records & 2,831 & 2,215 & 3,689 & $\mathbf{-21.8\%}$ & $\mathbf{+66.5\%}$ \\
pass yield (reads) & 348,399 & 387,614 & 366,625 & $+11.3\%$ & $-5.4\%$ \\
pass ambig codes & 1,615 & 533 & 927 & $\mathbf{-67.0\%}$ & $\mathbf{+73.9\%}$ \\
records w/ \texttt{err\_factor} $>$ 1 & 830 & 212 & 612 & $\mathbf{-74.5\%}$ & $\mathbf{+188.7\%}$ \\
yield w/ \texttt{err\_factor} $>$ 1 & 15,174 & 5,275 & 9,402 & $\mathbf{-65.2\%}$ & $\mathbf{+78.2\%}$ \\
\bottomrule
\end{tabular}
\end{center}

Reading the columns left-to-right: v0.7.7 (baseline) $\to$
\texttt{v08-strict} (new phases at 0.7-style permissivity) $\to$
v0.8.0 (also relaxes the defaults). The penultimate column shows the
new-phase contribution in isolation; the last column shows the
permissivity contribution. The two effects pull in opposite directions
on every metric.

\textbf{The new phases by themselves tighten the variant set in the
validator-preferred direction:} fewer pass-track records ($-22\%$), each
larger and better-supported ($+11\%$ yield per variant); ambiguity codes
slashed ($-67\%$); records with above-expected disagreement cut by
nearly three quarters ($-74\%$); the noise-blob, force-assigned, and
unphased-residual variants that 0.7 was emitting are largely caught and
routed elsewhere.

\textbf{The relaxed defaults push the variant count back up} ($+67\%$
records), pulling in many smaller candidate variants that the strict
thresholds would have rejected. Per-variant yield drops ($-5\%$),
ambiguity rises ($+74\%$), and records with above-expected disagreement
roughly triple ($+189\%$) as a consequence; the smaller candidate
variants are noisier than the larger anchors. Even after this
permissivity rebound, however, the v0.8 default remains better than
v0.7 on every quality axis: $-43\%$ ambig and $-26\%$
above-expected-disagreement records on a $+30\%$-larger pass set.

The $+30\%$ pass-record headline of v0.7 $\to$ v0.8 is the sum of these
two effects: $+67\%$ from looser thresholds, partially offset by $-22\%$
from the new phases doing their quality work.

\paragraph{Where the missing pass-yield reads go.} The
\texttt{v08-strict} $\to$ v0.8.0 pass-yield decrease ($-5.4\%$,
$-20{,}989$ reads) raised the question of whether reads were being lost.
They are not. Total \emph{core} yield in fact \emph{increases} $+3.4\%$
($+13{,}620$ reads) under v0.8 defaults; the apparent pass-track loss is
reads being routed to the \texttt{.ns} and \texttt{.lq} inspection
tracks rather than dropped:

\begin{center}
\begin{tabular}{lrrr}
\toprule
core $\to$ track & \texttt{v08-strict} & v0.8.0 & $\Delta$ reads \\
\midrule
core total & 396,385 & 410,005 & $+13{,}620$ \\
$\to$ pass & 387,614 & 366,625 & $-20{,}989$ \\
$\to$ \texttt{.ns} & 7,297 & 36,955 & $+29{,}658$ \\
$\to$ \texttt{.lq} & 1,474 & 6,425 & $+4{,}951$ \\
\bottomrule
\end{tabular}
\end{center}

The bookkeeping balances exactly: $-20{,}989 + 29{,}658 + 4{,}951 =
+13{,}620$. The mechanism is the chain we described in
Section~\ref{sec:results-cascade}: looser HP normalization
(\texttt{--hp-normalization-length 6} vs.~1) lets Phase~3 detect more
variant positions at HP boundaries and split clusters that
\texttt{v08-strict} kept whole. The smaller sub-clusters drop below
0.7-style thresholds but survive the v0.8 hard floor of 3 reads. They
then route to \texttt{.ns} (CER rejects them as plausible
basecaller-noise replicas of larger peers in the same identity group).
The pass-track decrease is the new significance filter doing exactly
the work it was designed for, on a larger candidate population that the
relaxed defaults made visible.

The takeaway for validators: the v0.7 $\to$ v0.8 ``$+30\%$ variants''
headline is permissivity-driven. The new phases by themselves produce
\emph{fewer}, \emph{better-supported}, \emph{less-ambiguous} pass-track
variants than v0.7 did. The variant-count gain you see in
v0.8-defaults is the new pipeline judging a larger candidate population
and surfacing more of it.

\paragraph{HP-interface variants surfaced.} Of the three permissivity
changes, the move from \texttt{--hp-normalization-length 1} (v0.7
default, all HP runs normalized) to \texttt{--hp-normalization-length
6} (v0.8 default, only runs $\geq$ 6 normalized) is the one specifically
motivated by the interface-homopolymer issue from
Section~\ref{sec:hp-and-ambiguity}. The validator-flagged class of
variants is structurally a 1-base shift in an HP run's boundary ---
e.g., \texttt{ACCTA} vs.~\texttt{ACTTA}, where the C/T interface has
moved by one position, with each side a homopolymer ``run'' (in
quotes because either side may be length 1, and the case generalizes
to longer runs). Under v0.7's HP-1 normalization, both forms compared
identically; the variant phaser never saw the interface position as
variable, reads from both alleles ended up in the same cluster, and
the resulting consensus either voted a majority HP-length (silently
discarding the minor allele) or emitted an ambiguity code at the
projected substitution position. Under v0.8's HP-6 normalization, HP
runs of length 1--5 are no longer normalized, so reads with different
HP-run lengths show real disagreement at the shift position; the
phaser detects this and splits the cluster.

To measure how much such heterogeneity v0.7 was hiding, we built a
phasing-implementation-independent detector that works directly on
cluster MSAs (Methods). For each cluster, the detector identifies every
consensus HP run of length $\geq 2$; for each read it scans left and
right of the consensus run, collapsing alignment gaps, to determine the
read's actual HP-run length at that position; the cluster is flagged
as having an unphased HP-interface if at least 10\% of reads (and at
least 3 reads, matching the speconsense phaser's variant-call
thresholds) show a single-direction $\pm 1$ shift. Restricted to HP
runs of length 2--5 (the population where v0.7 normalizes but v0.8
does not), and to pass-track core clusters only (the validator-facing
output):

\begin{center}
\begin{tabular}{lrr}
\toprule
metric & v0.7-defaults & v0.8-defaults \\
\midrule
pass-track core clusters analyzed & 2,209 & 3,034 \\
clusters with unphased HP-interface & 598 (27.1\%) & 102 (\textbf{3.4\%}) \\
pass-track specimens with $\geq$ 1 such cluster & 353 (49.0\%) & 91 (\textbf{11.9\%}) \\
\bottomrule
\end{tabular}
\end{center}

\textbf{Among pass-track clusters, the unphased-HP-interface rate drops
from 27\% in v0.7 to 3\% in v0.8 --- an 8-fold reduction --- and the
share of specimens whose pass-track output contains at least one such
cluster falls from 49\% to 12\%}, a 4-fold reduction. Most of v0.7's
heterogeneous clusters now phase into separate v0.8 sub-clusters with
their own consensuses, exposing the HP-interface variants v0.7 was
hiding. Two cross-checks support the detector: a structural enumeration
of within-specimen pass-track variant pairs (Methods) finds 7 pairs
across 5 specimens in v0.8 that differ only by a length-conserving
canonical HP-interface shift (the strict \texttt{ACCTA}
vs.~\texttt{ACTTA} case); v0.7 emits zero such pairs. Four of those
five specimens are independently flagged by the MSA-based detector
applied to v0.7's pass-track clusters; the fifth has a v0.7 cluster
showing HP-shift heterogeneity at sub-threshold magnitude (small
sample size, fewer than 3 reads in the shift direction), consistent
with the phaser's read-count minimum being the binding constraint.

The 91 v0.8 specimens with residual unphased HP-interface clusters are
mostly cases where the cluster's MSA is too small for the phaser's
read-count minimum, plus cases where the HP-shift co-occurs with
other variation that breaks the single-direction signal. The first is
fundamental; the second is plausibly addressable in a future revision.

\subsection{Cross-tool baseline: NGSpeciesID}
\label{sec:results-ngsi}

NGSpeciesID was developed for a different workload than speconsense
addresses. Sahlin and colleagues~\cite{ngspeciesid} optimized for
biodiversity-scale processing of single-amplicon specimens with minimal
preprocessing; isONclust~\cite{isonclust}, the clustering algorithm at
NGSpeciesID's core, uses minimizer hashing to skip pairwise alignment
for most reads, trading some clustering precision for the throughput
needed to process thousands of samples on a laptop. The design
assumption is one consensus per specimen --- the
\texttt{--abundance\_ratio} default of 10\% deliberately suppresses
minor variants --- and the authors explicitly note mixed-sample
handling as outside the design target.

The Russell protocol's invocation, refined over Mycota Lab's production
use, already opts out of much of that tradeoff:
\texttt{--mapped\_threshold 1.0} disables isONclust's minimizer-based
shortcut entirely (the threshold becomes unsatisfiable), forcing
parasail alignment for every clustering decision;
\texttt{--symmetric\_map\_align\_thresholds} requires the alignment
ratio to clear the threshold from both directions; and
\texttt{--abundance\_ratio 0.2} relaxes the variant-suppression default.
We ran NGSpeciesID 0.3.0 with medaka 2.0.0 polishing on the same
955-specimen ont98 dataset, with the full invocation \texttt{--ont
--consensus --symmetric\_map\_align\_thresholds --aligned\_threshold
0.75 --mapped\_threshold 1.0 --medaka --abundance\_ratio 0.2
--sample\_size 500 --top\_reads}. What we compare against speconsense
is therefore a precision-tuned NGSpeciesID pipeline --- greedy
clustering by parasail alignment, SPOA consensus, medaka polishing,
primer trim --- not isONclust-as-shipped.

NGSpeciesID numbers in the table below are unfiltered (every emitted
record, including RiC$<$3); the Mycota Lab production protocol applies
a RiC$\geq$3 post-filter to NGSpeciesID output before downstream use,
which §\ref{sec:results-tax} analyzes in parallel.

\begin{center}
\begin{tabular}{lrrr}
\toprule
metric & v0.7.7 & v0.8.0 & NGSpeciesID \\
\midrule
specimens with output & 837 & 850 & \textbf{926} \\
pass records & 2,831 & 3,689 & 1,180 \\
pass yield (reads) & 348,399 & 366,625 & 254,395 \\
pass ambig codes & 1,615 & 927 & \textbf{0} \\
records / specimen (mean) & 3.38 & 4.34 & 1.27 \\
records, \texttt{err\_factor} $>$ 1 & 830 (29\%) & 612 (17\%) & \textbf{431 (37\%)} \\
yield, \texttt{err\_factor} $>$ 1 & 15,174 (4\%) & 9,402 (3\%) & \textbf{78,019 (31\%)} \\
\bottomrule
\end{tabular}
\end{center}

NGSpeciesID is much more conservative on per-specimen variant discovery
(median 1 variant/specimen, mean 1.27) than either speconsense version.
This is by design: NGSpeciesID is built around the ``one consensus per
specimen'' use case, with its \texttt{--abundance\_ratio} filter cutting
clusters below 20\% of the dominant cluster's read count. The variants
speconsense surfaces beyond NGSpeciesID's count are the intragenomic
copies, paralogs, and minor allelic variation that motivated speconsense
in the first place.

NGSpeciesID's wider specimen coverage --- 89/76 more specimens with
output than speconsense 0.7/0.8 --- is partly a like-for-like artifact.
The counts above are unfiltered, and the Mycota Lab production protocol
post-filters NGSpeciesID's output to RiC$\geq$3 because single- and
two-read clusters are too uncertain to deposit (see
§\ref{sec:results-tax}). After that filter, NGSpeciesID and v0.8.0
cover essentially the same specimens (850 each, of 955). The
residual gap traces to speconsense's clustering splitting low-depth
read pools into sub-3-read clusters that fall below the emit
threshold, where NGSpeciesID's greedy single-cluster-per-specimen
design produces a consensus regardless. The trade-off is real ---
speconsense provides finer-grained variant resolution at the cost of
some specimens with very low input depth getting no output --- but it
is much narrower than the unfiltered numbers suggest, and what remains
is concentrated in exactly the single-read-anchored evidence the
protocol filter distrusts. Cross-referencing the 82 specimens that
NGSpeciesID's raw output covers but v0.8.0 does not against
§\ref{sec:results-tax}'s taxonomic agreement: 87\% of the 281 records
those specimens contribute are noise, only 10\% hit a primary target,
and 56 of the 82 specimens are truth-unknown to begin with. Most of
NGSpeciesID's wider raw coverage is noise emit from low-quality
specimens, and the protocol filter correctly rejects almost all of it.

The wider coverage comes with lower per-record quality, however. We
recomputed \texttt{err\_factor} for each NGSpeciesID record by
sampling 100 supporting reads, building an MSA via SPOA, and applying
speconsense's HP-aware comparator (the same code path the 0.8 pipeline
uses internally; see Methods). 37\% of NGSpeciesID's records show
\texttt{err\_factor} $>$ 1 and 31\% of its yield sits in those records
--- roughly \textbf{12$\times$ v0.8.0's $3\%$ yield share} above the
basecaller error model's expectation. \texttt{err\_factor} is a
model-relative metric computed under speconsense's HP-aware
comparator and Dorado-v5.0 \texttt{q\_ctx} tables; it measures
internal consistency under speconsense's error model rather than
absolute consensus quality. The records-vs.-yield ratio is
informative on its own: in v0.7/v0.8, above-expected-disagreement
records are concentrated in small clusters (records-share $\gg$
yield-share), but in NGSpeciesID the above-expected records are also
the large clusters (records-share $\approx$ yield-share). NGSpeciesID
has no equivalent of CER routing or \texttt{err\_factor} filtering;
records that speconsense would route to \texttt{.ns} or \texttt{.lq}
land directly in the pass output. A validator using NGSpeciesID would
need to apply downstream quality filtering to recover comparable
trustworthiness.

NGSpeciesID emits zero ambiguity codes by design: it has no mechanism
for surfacing residual within-cluster heterogeneity. A public reference
database may prefer the cleaner consensus; a validator looking for
unresolved heterogeneity loses information. Speconsense's
ambiguity-code emission, paired with the noise filter and the
significance framework, is one of the deliberate departures from the
NGSpeciesID model.

The differences in the table reflect different points in a design
space, not a defect on either side. NGSpeciesID was a good fit for
getting Mycota Lab's nanopore protocol off the ground, and remains a
sensible default for workflows whose primary target is one consensus
per specimen at biodiversity scale. The validator community's evolving
requirements --- finer variant resolution, defensible filtering of
low-support variants, ambiguity-code emission for residual
heterogeneity, and transparent rejection tracks --- pulled the workflow
toward a different optimization target.

\subsection{Taxonomic agreement}
\label{sec:results-tax}

The §\ref{sec:results-ngsi} comparison is at the level of records and
specimens --- counts and quality measures. The biology question is
whether those records correspond to real organisms. We added a
taxonomic-agreement step that classifies each emitted variant against a
per-specimen truth set: validator-confirmed organism IDs from the
iNaturalist observations behind 935 of the 955 ont98 specimens. The truth
set was built from a five-stage pipeline (vsearch + adjusted-identity
search against MycoMap, UNITE escalation for unknowns, an internal
reference of validator-confirmed (BISON, ITS) deposits, plus iNat
metadata and validator-comment review) and then manually walked, row by
row. After review, 857 specimens have a confirmed primary, 78 are
truth-unknown (specimen failures, contaminations, or validator-flagged
mixups), and 54 carry a recovered mycoparasite or co-occurring organism
alongside the primary. Methods §\ref{sec:methods-tax} describes the
pipeline. One framing note: the ont98 reference is built from
validator-confirmed deposits in the same dataset, so a specimen
matching its own ont98 deposit is a self-consistency check (the
pipeline's consensus reproduces the validated lab call from the same
reads) rather than a generalization test. We allow it because it is
the only way to score temp codes that haven't yet reached the public
reference --- most temp codes do appear in MycoMap, but a long tail of
recent or rarely-deposited ones do not. The MycoMap self-match filter
still applies for the specimens whose validator-confirmed call exists
in the public reference.

Each emitted variant lands in one of four buckets that sum to 100\% per
track. \textbf{primary\_target}: the variant canonically matches the
truth primary, or shares its genus (the same-genus extension absorbs
cases where the reference DB has only a genus-rank deposit, e.g.\
``Xylodon sp.'' covering ``Xylodon asper''). \textbf{myco}: a
parasite, host, or co-occurring organism on this specimen's substrate
(\emph{Trichaptum biforme} on a \emph{Phaeocalicium polyporaeum}
collection; \emph{Ampelomyces sp.} on \emph{Golovinomyces}; the
\emph{Syzygites megalocarpus} mycoparasite of an \emph{Agaricus} host
that itself failed to amplify). \textbf{contaminant}: known yeast or
upstream-process artifact. We treat cross-bleed contamination
(e.g.\ a high-RiC \emph{Hemileccinum hortonii} bleeding across multiple
unrelated specimens) as a contaminant rather than noise: it is an
artifact in the input data that any tool could in principle find, not
a tool-produced one. \textbf{noise}: everything else --- low-adj
non-matches, family-level fallbacks, unexplained sequences ---
potentially tool-produced.

Recall is defined as (specimens with at least one primary-target hit)
$/$ 955, computed against the full 955-specimen input
cohort.\footnote{The sequencing run loaded 960 specimens; 5 had no
reads after demultiplexing.} Three failure modes contribute to
non-recovery, and all count against recall: 20 of the 955 specimens
produced no record from any tool (low-yield wells with reads that fail
to cluster cohesively --- median 19 input reads each, mean read length
$\sim$1{,}300\,bp vs.~$\sim$700\,bp for the cohort, suggesting
concatemers or off-target amplification); 78 of the remaining
935 truth-set specimens are truth-unknown (specimen failures,
contaminations, validator-flagged mixups); and a tool may emit records
that simply do not match the truth primary. Pooling these on the same
denominator means coverage gaps surface as quality differences rather
than vanishing from the comparison.

\paragraph{NGSpeciesID's protocol filter splits the comparison.}
The Mycota Lab production protocol post-filters NGSpeciesID's output
to RiC$\geq$3, treating single- and two-read clusters as too uncertain
for downstream use. We apply that filter and treat the rejected
ric$<$3 emits as an \texttt{lq} track for parallel comparison with
v0.8's \texttt{.lq}. The case for the filter is two-fold. A single
read cannot be corrected for nanopore error, so even when the species
call is right, the sequence cannot be deposited in a reference
database --- the very pipeline that produces our internal ont98
reference enforces this constraint. And single-read cross-specimen
carry-over is common; we have at least one documented case (a
\emph{Hygrocybe}) where a validator initially selected a ric$=$1
sequence to identify the specimen, which in retrospect was a different
\emph{Hygrocybe} from another well's read pool, while a second ric$=$1
sequence in the same specimen actually matched the morphology. Our
20/20-hindsight truth set absorbs that particular swap, but the
protocol cannot afford to.

In the table below, the \texttt{ngsi-defaults pass} and \texttt{lq} rows
are the RiC$\geq$3 / RiC$<$3 partition of NGSpeciesID's output (873 +
307 = 1,180 unfiltered records, matching §\ref{sec:results-ngsi}). The
\texttt{ngsi-defaults lq} row's recall (2.3\%) and noise share (88.3\%)
look unhelpful in isolation; the paragraph ``The \texttt{.lq}-track
question'' below unpacks the 22 specimens that the RiC$\geq$3 filter
would lose were it not for those records.

\begin{center}
\begin{tabular}{lrrrrrrrr}
\toprule
run / track & cover & recall & records & \%prim & \%myco & \%cont & \%noise & lq\_resc \\
\midrule
ngsi-defaults pass     & 89.1\% & 83.7\% &   873 & 94.6\% & 0.7\% & 0.9\% &  3.8\% & 22 \\
ngsi-defaults lq       & 10.2\% &  2.3\% &   307 &  8.1\% & 1.9\% & 1.6\% & 88.3\% & --- \\
v07-defaults pass      & 87.6\% & 85.1\% & 2,831 & 85.5\% & 1.0\% & 2.5\% & 11.1\% &  0 \\
v08-strict pass        & 87.7\% & 85.3\% & 2,215 & 85.5\% & 1.1\% & 2.4\% & 10.9\% &  1 \\
\textbf{v08-defaults pass} & \textbf{89.0\%} & \textbf{86.3\%} & 3,689 & 82.1\% & 1.1\% & 2.4\% & 14.4\% &  0 \\
v08-defaults \texttt{.ns}  & 70.6\% & 69.0\% & 7,483 & 97.6\% & 0.4\% & 0.9\% & \textbf{1.0\%} & --- \\
v08-defaults \texttt{.lq}  & 52.5\% & 45.3\% & 1,659 & 89.3\% & 0.5\% & 2.0\% &  8.1\% & --- \\
\bottomrule
\end{tabular}
\end{center}

The recall ranking matches the coverage observations from
§\ref{sec:results-permissivity}/§\ref{sec:results-ngsi} but at the
cohort level. v0.8.0-defaults wins (86.3\%) by combining the highest
emit coverage (89.0\%) with comparable per-emitting recovery to the
strict variants. v0.7-defaults and v0.8-strict are essentially tied at
$\sim$85\% --- the +1.0-pt improvement v0.8-defaults shows over its
strict variant comes from the additional specimens its permissive
emit recovers, not from any in-emit recovery improvement. Permissivity,
not the new phases, is the recall lever.

The \%noise column is where the variant-detection vs.\ consensus design
difference shows up. NGSpeciesID is built around one record per
specimen, and 94.6\% of its records hit the primary target with only
3.8\% noise. v0.8-defaults emits 4.3 records per specimen on average
and 14.4\% of its pass records sit in the noise bucket. The naive read
of those numbers (``v0.8 is nearly four times noisier'') misses what
the records actually are, however: the noise concentrates dramatically
by identity group.

\paragraph{Noise concentrates in groups 2+.}
For the v0.7/v0.8 tools, an identity group is the cluster a record
came from --- group 1 is the rank-1-by-RiC cluster per specimen, where
the primary target is expected to live; groups 2+ are the additional
clusters surfaced by the variant-discovery pipeline.

\begin{center}
\begin{tabular}{lrrrr}
\toprule
run / track & group & records & \%prim & \%noise \\
\midrule
v07-defaults pass         & 1   & 2,438 & 96.9\% &  1.8\% \\
                          & 2+  &   393 & 14.2\% & 68.4\% \\
v08-strict pass           & 1   & 2,079 & 90.0\% &  7.5\% \\
                          & 2+  &   136 & 16.9\% & 63.2\% \\
v08-defaults pass         & 1   & 3,040 & \textbf{97.2\%} &  \textbf{1.8\%} \\
                          & 2+  &   649 & 11.4\% & 73.5\% \\
v08-defaults \texttt{.ns} & 1   & 7,361 & \textbf{98.9\%} &  \textbf{0.2\%} \\
                          & 2+  &   122 & 23.8\% & 52.5\% \\
v08-defaults \texttt{.lq} & 1   & 1,518 & 96.6\% &  2.8\% \\
                          & 2+  &   141 & 11.3\% & 65.2\% \\
\bottomrule
\end{tabular}
\end{center}

A validator inspecting only the rank-1 cluster sees noise drop from
14\% to 2\% in v0.8.0's pass track, and from 1\% to 0.2\% in the
\texttt{.ns} track. The group-2+ records do carry real biology ---
11--17\% are primary-target hits and an additional 4--7\% are
myco-associates or known contaminants --- but the majority are
cross-bleed or unexplained. NGSpeciesID emits at most one consensus
per specimen-cluster (no phasing within), so this analysis is specific
to the variant-discovery tools. The validator-facing implication is
that group 1 alone covers nearly all the recall while contributing a
small fraction of the noise; groups 2+ are worth scanning for
co-occurring biology but are noisier per record.

\paragraph{The \texttt{.lq}-track question.}
Both speconsense \texttt{.lq} and NGSpeciesID-protocol-filtered ric$<$3
represent records the tool emitted but the protocol/routing flagged as
low quality. Do those records ever recover specimens the pass track
misses? For NGSpeciesID, 22 specimens are recovered \emph{only} by
ric$<$3 emits --- specimens whose only sequence evidence is a
single-read or two-read cluster. The ric$\geq$3 filter rejects these
entirely; without the \texttt{lq} treatment they would be invisible
recall failures. Of the 7.9-percentage-point coverage drop between
NGSpeciesID's raw (97.0\%) and protocol-filtered (89.1\%) outputs, 22
specimens correspond to validator-confirmed primaries the protocol
filter loses. The other side of that ledger: the 75 specimens lost to
the filter overall contribute 264 records, of which 87.5\% are noise
--- the filter rejects 22 valid recoveries while also rejecting
$\sim$240 noise records. The two-fold filter justification still holds,
and the trade favors the filter, but the recall cost is real and
visible.

For v0.8.0, 0 specimens are rescued by \texttt{.lq}: every specimen
the \texttt{.lq} track recovers is already covered by pass.
\texttt{.lq} is purely a within-specimen variant-exploration track ---
it surfaces additional clusters for already-recovered specimens, not
new specimen coverage. The asymmetry highlights how each tool handles
low-quality clusters: NGSpeciesID's ric$<$3 emits include a tail of
legitimate single-read specimens that protocol post-filtering loses,
while v0.8.0's MCL clustering and routing capture the same specimens
through normal pass channels.

\paragraph{Mycoparasites and substrate co-occurrences.}
54 specimens have a recovered organism alongside the primary that the
truth set classifies as a mycoparasite, host, or substrate
co-occurrence. The cleanest examples are the obligate parasitic
relationships: \emph{Phaeocalicium polyporaeum} growing on
\emph{Trichaptum biforme} (the validator's iNat note recorded only the
\emph{Trichaptum} sequence, but v0.8.0 cleanly phased both the
Phaeocalicium target and the Trichaptum host); \emph{Ampelomyces sp.}
on a \emph{Golovinomyces} powdery mildew (the validator separately
created a duplicate observation for the mycoparasite); and
\emph{Syzygites megalocarpus} on an \emph{Agaricus} that itself failed
to amplify but where the parasite came through cleanly. Recall on the
mycoparasite sub-cohort tracks the per-record \%myco column: NGSpeciesID's
abundance-ratio default suppresses minor variants and rarely surfaces
co-occurring organisms (0.7\% myco), while v0.8.0's variant phasing
preserves them at the same rate as v0.7 ($\sim$1.1\%) but on a larger
record base.

\section{Methods}
\label{sec:methods}

\textbf{Dataset.} ont98: 955 demultiplexed per-specimen FASTQ files from a
single Mycota Lab production run, Russell ITS amplicon
protocol~\cite{russell2023}, Dorado v5.0 SUP basecalled. Same primer set
(ITS1F / ITS4) across all specimens. Per-specimen read counts range from a
few hundred to a few thousand.

\textbf{Cross-version and MAD-ablation runs.} Two comparisons drawn
from fresh core runs underlie the body's quantitative claims outside
the additive ladder:

\begin{center}
\begin{tabular}{l l p{0.42\linewidth}}
\toprule
Comparison & Reference & Run \\
\midrule
Cross-version (Sections~\ref{sec:results-version}--\ref{sec:results-tax})
  & v0.7.7 defaults & v0.8.0 defaults \\
Phase 9 MAD (Sections~\ref{sec:results-additive}, \ref{sec:results-routing})
  & v0.8 defaults & \texttt{--disable-mad-outlier-removal} (\texttt{v08-no-mad}) \\
\bottomrule
\end{tabular}
\end{center}

The \texttt{--disable-mad-outlier-removal} flag was added to the
speconsense codebase specifically to support this study; it ships as
part of 0.8.0. Per-phase subtractive ablations for the other new phases
were also run during pipeline development but are superseded by the
additive ladder below for the body's analyses.

\textbf{Permissivity-isolation run.} One additional fresh core run,
\texttt{v08-strict}: the v0.8.0 binary with v0.7-style restrictive
defaults (\texttt{--min-size 5 --min-cluster-ratio 0.01
--hp-normalization-length 1}) and otherwise default summarize. Used in
Section~\ref{sec:results-permissivity} to decompose the v0.7~$\to$~v0.8
delta into new-phase and permissivity components.

\textbf{Cross-tool baseline.} NGSpeciesID 0.3.0 with medaka 2.0.0
polishing, run via the local
checkout\footnote{\url{https://github.com/joshuaowalker/NGSpeciesID} ---
the brew/PyPI 0.3.0 release does not include the
\texttt{--symmetric\_map\_align\_thresholds} or \texttt{--top\_reads}
flags used in the Russell protocol.} at \texttt{\textasciitilde
/mm/code/NGSpeciesID} with the Russell-protocol invocation:
\texttt{--t 1 --ont --consensus --symmetric\_map\_align\_thresholds
--aligned\_threshold 0.75 --mapped\_threshold 1.0 --medaka
--abundance\_ratio 0.2 --sample\_size 500 --top\_reads --primer\_file
\$PRIMERS}. Output collected into mycomap-format summary via the same
\texttt{summarize.py} that Mycota Lab uses in production. Used in
Section~\ref{sec:results-ngsi}.

\textbf{Additive-ladder matrix.} Six fresh core runs on the v0.8.0
binary, starting from a stripped-down baseline and adding one new phase
at a time:

\begin{center}
\begin{tabular}{l p{0.55\linewidth}}
\toprule
Run ID & Disable flags (everything else default) \\
\midrule
\texttt{v08-add-base} & \texttt{--disable-noise-filter --disable-read-reassignment --disable-mad-outlier-removal} \\
\texttt{v08-add-p6} & \texttt{--disable-noise-filter --disable-discard-recovery --disable-second-phasing --disable-mad-outlier-removal} \\
\texttt{v08-add-p6-p5} & \texttt{--disable-discard-recovery --disable-second-phasing --disable-mad-outlier-removal} \\
\texttt{v08-add-p6-p7} & \texttt{--disable-noise-filter --disable-second-phasing --disable-mad-outlier-removal} (alt-order) \\
\texttt{v08-add-p6-p5-p7} & \texttt{--disable-second-phasing --disable-mad-outlier-removal} \\
\texttt{v08-add-p6-p5-p7-p8} & \texttt{--disable-mad-outlier-removal} \\
\bottomrule
\end{tabular}
\end{center}

The terminal step (full v0.8.0 with all phases enabled, including MAD
and CER/err\_factor routing) is the existing \texttt{v08-defaults} run
from the subtractive matrix, reused. Each additive run is also
re-summarized in \emph{permissive} mode (\texttt{--min-cer-factor=0
--max-err-factor=0}) so that core conservation can be verified
unambiguously. Phases~10 and~12's CER/err-factor routing are summarize-only
and apply uniformly across the additive ladder via the default
summarize step.

\textbf{Cross-version \texttt{err\_factor}.} For v0.7 runs (where
\texttt{err\_factor} is not present in FASTA header output), values were
computed offline from the per-cluster \texttt{cluster\_debug/*-msa.fasta}
files using the canonical 0.8 implementation. Cross-validation against the
FASTA-header values for v0.8 runs showed agreement to within floating-point
rounding (max abs diff 0.0005 across all 13,300 v0.8 core records).

\textbf{NGSpeciesID \texttt{err\_factor}.} NGSpeciesID does not compute
\texttt{err\_factor}, so values were derived externally for the
cross-tool comparison. For each NGSpeciesID record we sampled up to 100
supporting reads from
\texttt{clusters/<specimen>/reads\_to\_consensus\_<id>.fastq} (matching
speconsense's production sample size), built an MSA via SPOA with the
same options speconsense uses internally
(\texttt{-r 2 -l 1 -m 1 -n -1 -g -1 -e -1}), and applied
\texttt{compute\_cluster\_err\_factor} from \texttt{speconsense.msa} ---
the canonical 0.8 implementation, with the same
\texttt{dorado-v5.0} error-model table. The MSA reference is SPOA's own
consensus from the sampled reads, so the value reported is the
per-cluster read-to-consensus disagreement ratio that
\texttt{err\_factor} measures throughout this paper.

\textbf{HP-interface analysis.} Two complementary detectors targeted
the interface-homopolymer effect. (1) \emph{Structural pair
enumeration.} For each run we enumerated every within-specimen pair
of pass-track records and classified them structurally: each sequence
was decomposed into a list of (base, run-length) tuples, and pairs
were classified by relating the two decompositions. A pair was tagged
\texttt{hp\_interface\_1\_5} if the two decompositions had identical
shape (same number of runs, same base at each position), all
length-differing runs were $\leq 5$, and the differences came in
adjacent $+N/-N$ pairs (the canonical length-conserving boundary shift
\texttt{ACCTA}~vs.~\texttt{ACTTA}); \texttt{hp\_extension\_1\_5} if
shape matched and all differing runs were $\leq 5$ but the differences
were not strictly adjacent $\pm N$ pairs (single-run length changes,
non-adjacent shifts). Implementation in
\texttt{validation/scripts/hp\_interface\_pairs.py}; per-pair output
in \texttt{validation/analysis/hp-interface-pairs.csv}.
(2) \emph{MSA-level unphased-interface detection.} The structural
enumeration is conservative: it only catches pairs whose differences
are exclusively HP-shift, missing real HP-interface variants that
co-occur with other variation. The MSA-level detector, applied to
each cluster's \texttt{cluster\_debug/<id>-msa.fasta}, addresses this.
For each consensus HP run of length $\geq 2$, the detector computes
each read's actual HP-run length at that position by anchoring at
the consensus run's MSA columns, walking outward in the read while
collapsing alignment gaps, and stopping at the first non-gap base
that is not the run's character. The cluster is flagged as having an
unphased HP-interface at that run if at least 10\% of reads, and at
least 3 reads, show a single-direction $\pm 1$ shift relative to the
consensus's run length --- matching the speconsense phaser's
variant-call thresholds. Reads with multi-base shifts ($\pm 2$ or
more) and reads with substitutions inside the run are excluded from
the shift count to avoid conflating HP-interface signal with other
variation. Implementation in
\texttt{validation/scripts/hp\_interface\_msa.py}; per-flagged-run
output in
\texttt{validation/analysis/hp-interface-msa-v07.csv} and
\texttt{-v08.csv}.

\textbf{Taxonomic agreement.} \label{sec:methods-tax} The
§\ref{sec:results-tax} truth set was constructed in five stages.
(1)~A union FASTA was built by deduplicating every consensus emitted
across all run/track combinations by sequence SHA-1, yielding $\sim$30k
unique sequences against the four-tool $\times$ all-track product.
(2)~Each unique sequence was searched against three references in
priority order: the MycoMap fungal-ITS reference (103,465 records),
UNITE 19.02.2025 (2.07M records, used for unknowns at MycoMap-best
adj$<$0.95), and an internal ont98 reference of 576 records built from
iNaturalist observations where both the validator-confirmed BISON
(``Barcode Inferred Species or Name'') and the DNA-Barcode-ITS field
were populated, plus two manual additions for specimens where the
validator confirmed a species in comments without setting BISON. The
search uses vsearch \texttt{usearch\_global} as a fast filter followed
by a per-hit pairwise rescoring with the HP- and IUPAC-aware
\texttt{adjusted-identity} comparator (the same comparator the 0.8
pipeline uses internally). The MycoMap self-match filter is keyed by
iNat ID embedded in MycoMap accessions: a specimen's own MycoMap
deposit is excluded from its candidate set so that recall against an
independent public reference is not measured trivially. The ont98
internal reference is deliberately exempt from this filter --- its
records are validator-confirmed (BISON, ITS) pairs from the same
dataset, and self-matching is the mechanism by which we score
validator-confirmed temp codes that have not yet reached MycoMap ---
most temp codes do appear in the public reference, but a long tail of
recent or rarely-deposited ones do not.
Matching a specimen against its own ont98 deposit is therefore a
self-consistency check (does the pipeline's consensus reproduce the
validated lab call from the same reads?) rather than a generalization
test against an independent reference. We treat it that way in the
results.
(3)~iNaturalist metadata was pulled for each specimen via the v1 API
--- community taxon, OFV (observation field) values, comments, and
identifications --- with comments and identifications tagged by
validator role from a manually-curated rank list of validator and
expert iNat usernames. (4)~A draft truth file was synthesized by
combining BISON, PSN, the iNat species/genus, the search hits, and
validator commentary, with thirteen boolean flags surfacing rows
needing human review (genus-only matches, morphology-only validation,
high-RiC non-primary candidates, and so on). (5)~The 935-row draft
was manually walked by the author. Review was \emph{not} blinded to
tool output --- the draft surfaces per-tool routing decisions
alongside iNat metadata and sequence-search hits, and the author had
context on which records came from which run --- so the truth set
should be read as an author-curated reference standard, not a
blinded ground truth. Adjudication priority was: validator-confirmed
BISON or DNA-Barcode-ITS values from iNat; iNat community-taxon
species rank with sequence agreement; validator commentary on the
specimen; and biological plausibility against MycoMap, UNITE, and
ont98 search hits. $\sim$574 trivial cases (BISON-validated or iNat
species match) were spot-checked and bulk-accepted; the remaining
$\sim$361 decision rows were resolved one by one.
The final truth set assigns each specimen a confirmed primary taxon (or
truth-unknown), a list of mycoparasites/co-occurring organisms, and a
list of contaminants, plus a confidence tag distinguishing
high-confidence calls (validator-confirmed via BISON or DNA-Barcode-ITS)
from low-confidence calls (e.g.\ NGSpeciesID-only ric$<$3 emits where
the truth assignment depends on the very record being scored). All
truth-construction scripts and the final truth.csv are committed under
\texttt{validation/scripts/} and
\texttt{validation/analysis/taxonomy/}.

\textbf{Reproducibility.} All metric extraction, aggregation, and pairwise
comparison scripts, along with the per-run aggregate CSVs and per-hypothesis
comparison tables, are committed under \texttt{validation/} in the paper
repository.

\section{Looking ahead}
\label{sec:looking-ahead}

The 0.8.0 architecture is built around four ideas: annotate everything,
decide downstream; separate noise filtering from significance testing;
phase variants aggressively to reduce ambiguity codes; and design the
new phases as a coupled cascade rather than a list of independent
options. The first makes the pipeline replayable --- a validator who
disagrees with a default threshold can change it without re-running
the (expensive) clustering and consensus steps. The second makes the
validator's mental model match the pipeline: \texttt{err\_factor} for
``is this cluster real?'' and \texttt{cer\_factor} for ``is this
variant distinguishable from its peers?'' are the two questions a
validator naturally asks, and the pipeline now produces a number for
each. The third inverts the older default: where 0.7 reached for an
ambiguity code whenever phasing was inconclusive, 0.8 makes a
sustained effort to phase the variation cleanly through Phase~3, the
noise-filter-mediated cleanup of Phases~5--7, and the second-pass
phasing of Phase~8 --- and only emits an ambiguity code when none of
those mechanisms can resolve the disagreement.

The empirical results above support all four concretely. The CER
framework's \texttt{.ns} track contains variants whose \texttt{cer\_factor}
median is 0.022 --- not borderline calls but variants the math identifies
with high confidence as basecaller-noise replicas. Validators previously
dismissed these by hand; the pipeline now does the routing for them.
Phase~9's MAD outlier removal compresses the \texttt{err\_factor}
distribution enough to make \texttt{--max-err-factor=1.5} a defensible
default (§\ref{sec:results-routing}). Discard recovery
(Phase~7), once Phase~5 has filtered the discard pool, contributes a
$+18\%$ pass-track yield gain. Every core record is accounted for in
pass, \texttt{.ns}, or \texttt{.lq} (with summarize identity-group
merging conflating some pass-track variants), making every routing
decision auditable.

The fourth idea was less visible in the design but is visible in the
data: the new phases were designed as a coupled cascade, and the
additive matrix shows that they are. Adding any single phase in isolation can move
results in unexpected directions ---
notably, adding discard recovery (Phase~7) without first adding the
noise filter (Phase~5) actively increases pass-track ambiguity, because
unvetted reads from noise-blob clusters get pulled back into otherwise-clean
ones. The intended benefit only appears when the {P5, P6, P7, P8} group
is enabled together. The ablation flags shipped with 0.8.0 are intended
for studies of this kind, not as production knobs --- a user who wants
to keep one of these phases off in production is, on the evidence, asking
for a different tool.

What remains for future work breaks into two horizons: a near-term
0.8.1 fast-follow that ships parameter-tuning refinements identified
during the 0.8.0 validation, and a 0.8.2 revision focused on the
defaults that were left as deliberate first cuts pending
production-cohort evidence.

The 0.8.1 fast-follow ships a tuned MAD outlier-removal default. A
sweep of the four MAD CLI parameters across the full 955-specimen
cohort found that the four knobs split into two functional families
with opposing effects on \texttt{err\_factor}'s discriminative power.
The shape-aware rules (\texttt{--mad-z-threshold},
\texttt{--mad-gap-factor}) flag statistical outliers within a cluster
regardless of magnitude --- selectively present in primary clusters
whose heterogeneity comes from sequencing noise around a coherent
consensus, and largely absent from noise clusters whose heterogeneity
is distributed across all reads. Tightening these rules tightens
primary-cluster \texttt{err\_factor} distributions without disturbing
noise clusters, increasing the separation between the two populations
and making any choice of \texttt{--max-err-factor} a more accurate
classifier. The magnitude-aware safety floor
\texttt{--mad-min-drop-from-median} does the opposite: it compresses
\texttt{err\_factor} uniformly across primary and noise clusters,
decreasing the routing classifier's accuracy even though individual
cluster consensuses look cleaner. The fourth knob
(\texttt{--mad-min-mad}) had no detectable effect at any tested
setting. On the basis of the sweep, 0.8.1 tightens
\texttt{--mad-z-threshold} from 3.0 to 1.5: the AUC of
\texttt{err\_factor} as a noise-vs-primary classifier rises from 0.63
to 0.66, the F1 of the \texttt{--max-err-factor=1.5} routing decision
rises by 30\%, and per-specimen primary recall moves by one specimen
(within slack).

The 0.8.2 revision targets the routing and size-based defaults left
as first cuts in 0.8.0. \texttt{--min-cer-factor=1.0} is theoretically
grounded as the basecaller-noise line and is unlikely to move.
\texttt{--max-err-factor=1.5} and the size-based thresholds
(\texttt{--min-size}, \texttt{--min-cluster-ratio}) are safe but
unoptimized; the \texttt{v08-strict}/\texttt{v08-defaults} contrast in
§\ref{sec:results-permissivity} highlights the size-threshold
question in particular, where the right operating point sits
somewhere between the two. A joint sweep of \texttt{--max-err-factor}
under the 0.8.1 MAD configuration found that tightening to 1.0 or 1.2
costs 6--12 primary recoveries and is not currently defensible; size
thresholds remain on the watch list pending production-cohort
evidence. We expect 0.8.2 to settle those defaults based on
key-partner review, validator feedback, and user-acceptance testing
on additional production cohorts.

The longer horizon includes refreshing the error-model tables as
basecallers evolve, validating the CER framework against orthogonal
datasets, and exploring whether the framework extends gracefully to
non-amplicon use cases. The most open challenge is downstream:
choosing how to surface the multi-track output (pass / \texttt{.ns} /
\texttt{.lq}) to validators and to public reference databases in a way
that preserves the information content the pipeline has worked to
surface, without overwhelming downstream consumers who are used to one
consensus per specimen. Work the pipeline has done to expose
biological variation cleanly is only useful if the consumers of its
output can act on it.

\section*{Acknowledgements}

The CER framework grew out of a sustained conversation between Stephen
Russell (Mycota Lab), Harte Singer (Fungal Diversity Survey), and the
author, all three concerned about where to draw the line between a real
low-support variant and basecaller noise; specific cluster-level questions
Stephen flagged about specific specimens, and the SNP-calling-at-scale
problem Harte raised in August 2025, motivated the work directly. Two
sequence validators in particular drove the recurring questions of
Section~\ref{sec:hp-and-ambiguity} into the empirical work behind 0.8.0:
Elora Adamson, Local Coordinator for the MycoMap BC Network project
(British Columbia), surfaced both the interface-homopolymer issue that
drove the context-aware error model and the ambiguity-inflation issue
that drove the noise filter, with a running set of cluster-level
examples on each. Jessica Williams, of the Volunteer Sequence Validation
Corps (VSVC), provided detailed reviews of pre-0.7 outputs through
January 2026 that surfaced the merge-driven ambiguity-injection bug
fixed by raising \texttt{--merge-min-size-ratio} in v0.7, the
homopolymer-interface question, and a partial-overlap chimera that
informed the v0.6.1 fix. We are also grateful to the broader community
of sequence validators across the fungal-barcoding world: every
recurring question in this paper traces back to a real specimen
surfaced by a real validator, and the 0.8.0 pipeline is in large
measure a response to their patience.

All sequencing data analyzed in this paper were produced by Mycota Lab
(Hoosier Mushroom Society) under the protocol of
Russell~\cite{russell2023}. The author thanks Stephen D.~Russell for
sharing the run-level outputs that made the empirical sections of this
paper possible.

Substantial portions of the text were drafted with the assistance of
Claude (Anthropic, Opus~4.7). The core ideas, design decisions, and
empirical analyses were provided by the author; the AI contributed to
exposition, \LaTeX{} formatting, and iterative drafting under human
direction. The author takes full responsibility for the content.

\bibliographystyle{plainnat}
\begin{thebibliography}{99}

\bibitem[Walker, 2026a]{cer_paper}
J.~Walker.
\newblock Critical error rate analysis for {MSA}-based variant phasing in
nanopore amplicon sequencing.
\newblock Technical report, 2026.
\newblock \url{https://github.com/joshuaowalker/speconsense/blob/main/docs/variant_significance/variant_significance.pdf}

\bibitem[Walker, 2026b]{hp_paper}
J.~Walker.
\newblock Homopolymer error rate analysis for nanopore amplicon sequencing.
\newblock Technical report, 2026.
\newblock \url{https://github.com/joshuaowalker/speconsense/blob/main/docs/hp_error_rate/hp_error_rate_report.pdf}

\bibitem[Walker, 2026c]{cer_in_practice}
J.~Walker.
\newblock Critical error rate in practice: context-aware variant significance
for nanopore amplicon pipelines.
\newblock Technical report, 2026.
\newblock \url{https://github.com/joshuaowalker/speconsense/blob/main/docs/cer_in_practice/cer_in_practice.pdf}

\bibitem[Walker, 2024]{speconsense}
J.~Walker.
\newblock Speconsense: high-quality clustering and consensus generation for
{Oxford Nanopore} amplicon reads.
\newblock \url{https://github.com/joshuaowalker/speconsense}, 2024.

\bibitem[Russell, 2023]{russell2023}
S.~D. Russell.
\newblock {ONT DNA} barcoding fungal amplicons w/ {MinION \& Flongle V.3}.
\newblock \emph{protocols.io}, March 2023.
\newblock \url{https://dx.doi.org/10.17504/protocols.io.36wgq7qykvk5/v3}.

\bibitem[Sahlin et~al., 2021]{ngspeciesid}
K.~Sahlin, M.~Lim, and S.~Prost.
\newblock {NGSpeciesID}: {DNA} barcode and amplicon consensus generation from
long-read sequencing data.
\newblock \emph{Ecology and Evolution}, 11(3):1392--1398, 2021.

\bibitem[Sahlin and Medvedev, 2019]{isonclust}
K.~Sahlin and P.~Medvedev.
\newblock De~Novo clustering of long-read transcriptome data using a greedy,
quality-value based algorithm.
\newblock In \emph{Research in Computational Molecular Biology (RECOMB 2019)},
Lecture Notes in Computer Science, pages 227--242. Springer, 2019.

\bibitem[van~Dongen, 2000]{vandongen2000}
S.~van~Dongen.
\newblock \emph{Graph Clustering by Flow Simulation}.
\newblock PhD thesis, University of Utrecht, 2000.

\bibitem[Lee, 2003]{lee2003}
C.~Lee.
\newblock Generating consensus sequences from partial order multiple sequence
alignment graphs.
\newblock \emph{Bioinformatics}, 19(8):999--1008, 2003.

\bibitem[Vaser et~al., 2017]{vaser2017}
R.~Vaser, I.~Sovi\'{c}, N.~Nagarajan, and M.~\v{S}iki\'{c}.
\newblock Fast and accurate de novo genome assembly from long uncorrected
reads.
\newblock \emph{Genome Research}, 27(5):737--746, 2017.

\bibitem[\v{S}o\v{s}i\'{c} and \v{S}iki\'{c}, 2017]{edlib}
M.~\v{S}o\v{s}i\'{c} and M.~\v{S}iki\'{c}.
\newblock Edlib: a {C/C++} library for fast, exact sequence alignment using
edit distance.
\newblock \emph{Bioinformatics}, 33(9):1394--1395, 2017.

\bibitem[Rognes et~al., 2016]{vsearch}
T.~Rognes, T.~Flouri, B.~Nichols, C.~Quince, and F.~Mah\'{e}.
\newblock {VSEARCH}: a versatile open source tool for metagenomics.
\newblock \emph{PeerJ}, 4:e2584, 2016.

\bibitem[Walker, 2025]{adjusted_identity}
J.~Walker.
\newblock \texttt{adjusted-identity}: a {Python} package for
homopolymer-aware sequence identity scoring.
\newblock \url{https://github.com/joshuaowalker/adjusted-identity}, 2025.

\end{thebibliography}

\end{document}
